% Companion workbook for the IAA-SO School on AI/ML in Astronomy 2026
%   Unsupervised Learning pillar: chemical tagging of star clusters.
%
% Built on the Annual Reviews ar-1col class (vendored, unmodified) with the
% publisher's yoke replaced by the school's own information; see README.md.
%
% Build:  make            (vendors the deck figures, then latexmk)
%         make pdf        (LaTeX only)
\documentclass{ar-1col}

\setcounter{secnumdepth}{3}

\usepackage{url}
\usepackage{graphicx}
\usepackage{amsmath}
\usepackage{amssymb}
\usepackage{booktabs}
\usepackage{longtable}
\usepackage{array}
\usepackage{textcomp}
\usepackage[T1]{fontenc}
\usepackage[authoryear,round]{natbib}
\bibpunct{(}{)}{;}{a}{}{,}

% -------------------------------------------------------------------- links
% Clickable URLs. hyperref has to be loaded after every other package, which is
% why it sits here rather than beside \usepackage{url} above. hidelinks keeps
% the printed page exactly as it was: the URL text already reads as a URL, so a
% coloured box would only fight the Annual Reviews layout. breaklinks lets a
% long URL wrap inside a caption instead of running past the right margin.
\usepackage[hidelinks,breaklinks]{hyperref}

% The four addresses the workbook points at, as clickable \href targets. The
% display forms carry \allowbreak so a long URL wraps at its slashes, and they
% match the plain-text forms printed on the title page (\@deckurl and friends).
%
% KEEP THESE IN STEP with scripts/sync_figures.py, which encodes the same four
% URLs into the QR PNGs. A caption link that disagrees with the QR printed
% beside it is worse than no link at all.
\newcommand{\wbDeckURL}{garciadias.github.io/\#/\allowbreak presentations/\allowbreak iaa-so-chemical-tagging-2026}
\newcommand{\wbDeckLink}{https://garciadias.github.io/\#/presentations/iaa-so-chemical-tagging-2026}
\newcommand{\wbRepoURL}{github.com/\allowbreak garciadias/\allowbreak iaa-advanced-neural-networks-2026-draft}
\newcommand{\wbRepoLink}{https://github.com/garciadias/iaa-advanced-neural-networks-2026-draft}
\newcommand{\wbSchoolURL}{granadacongresos.com/ai-ml}
\newcommand{\wbSchoolLink}{https://www.granadacongresos.com/ai-ml}
\newcommand{\wbLibraryURL}{github.com/\allowbreak garciadias/\allowbreak iaa-advanced-neural-networks-2026-draft/tree/main/\allowbreak docs}
\newcommand{\wbLibraryLink}{https://github.com/garciadias/iaa-advanced-neural-networks-2026-draft/tree/main/docs}

% ------------------------------------------------------------------ metadata
\jname{IAA-SO School on AI/ML in Astronomy 2026 · companion workbook}
\jvol{Unsupervised Learning}
\jyear{2026}
\doi{}

% The Annual Reviews yoke (journal, volume, doi, copyright line) is replaced
% by the school's own information, per the brief.
\makeatletter
\def\@deckurl{garciadias.github.io/\#/presentations/\allowbreak iaa-so-chemical-tagging-2026}
\def\@repourl{github.com/\allowbreak garciadias/\allowbreak iaa-advanced-neural-networks-2026-draft}
\def\@schoolurl{granadacongresos.com/ai-ml}
\def\journalinfo{{\jinfofont
  Companion workbook for the IAA-SO School on AI/ML in Astronomy 2026
  (Unsupervised Learning pillar)\endgraf
  \vskip4\p@ School programme: \texttt{\@schoolurl}\endgraf
  \vskip4\p@ Live deck (animated figures): \texttt{\@deckurl}\endgraf
  \vskip4\p@ Code, data, results: \texttt{\@repourl}\endgraf
  \vskip4\p@ \copyright\ \@jyear\ Rafael Garcia-Dias. Prepared for teaching use.\endgraf}}
% The class sets the Annual Reviews offprint trim (7.375 x 9.125 in), but its
% own title-page block is 533 pt wide, so the offprint page clips it under
% pdfTeX. Set the page to letter, which the layout was cut for (456 pt text
% block plus margin-note gutter), and pin the PDF page size to match.
\setlength{\paperwidth}{8.5in}
\setlength{\paperheight}{11in}
\setlength{\pdfpagewidth}{\paperwidth}
\setlength{\pdfpageheight}{\paperheight}
% the class prints www.annualreviews.org in the running footer; replace it with
% the school's own line. NB the class already called \pagestyle{headings} while
% loading, which *executed* its footer definition -- so the redefinition below
% must be followed by another \pagestyle{headings} to take effect.
\def\ps@headings{%
  \let\@mkboth\@gobbletwo
  \def\@oddfoot{\hbox to\textwidth{\rhfont IAA-SO 2026 · chemical tagging\hfill\foliofont\thepage}}%
  \def\@evenfoot{\hbox to\textwidth{\foliofont\thepage\hfill\rhfont Companion workbook: chemical tagging}}%
  \def\@evenhead{}%
  \def\@oddhead{}%
  \def\sectionmark##1{}%
  \def\subsectionmark##1{}}%
\makeatother

\pagestyle{headings}

% The class typesets the contents as a single unbreakable shaded \vbox, 36pc
% wide and shifted 3.75pc sideways -- designed for a journal's handful of
% sections. A 17-chapter workbook overflows it by 490pt and bleeds off the
% left edge of the page, so the contents is re-set here as ordinary,
% page-breakable material.
\makeatletter
\renewcommand\tableofcontents{%
  \section*{\contentsname}%
  \@starttoc{toc}}
% ...and its contents-line number boxes are 1.2em wide, which two-digit
% chapter numbers fill completely, welding "10.3." to the title. Widen them.
\renewcommand*\l@section{\@dottedtocline{1}{0em}{2.4em}}
\renewcommand*\l@subsection{\@dottedtocline{2}{2.4em}{3.4em}}
\makeatother

% ------------------------------------------------------------------- helpers
% The ar-1col class recomputes every figure at its *natural* size: its
% \Gin@setfile hook rebuilds the box at \Gin@req@width and merely centres it,
% so the width= key is ignored and a 1456-px PNG lands 372 pt wide whatever
% you ask for. Everything is therefore wrapped in \resizebox, which acts on the
% box after the class has had its way.
%   \fig[<fraction of \textwidth>]{<file in figures/>}
\newcommand{\fig}[2][0.94]{\resizebox{#1\textwidth}{!}{\includegraphics{figures/#2}}}
% One figure, captioned, from figures/<file>.
\newcommand{\wbfig}[4][0.86]{%
  \begin{figure}[htbp]\centering
    \includegraphics[width=#1\textwidth]{figures/#2}
    \caption{#3}\label{#4}
  \end{figure}}
% Two panels side by side.
\newcommand{\wbtwo}[4]{%
  \begin{figure}[htbp]\centering
    \begin{tabular}{@{}c@{\hskip4pt}c@{}}
      \fig[0.46]{#1} &
      \fig[0.46]{#2}
    \end{tabular}
    \caption{#3}\label{#4}
  \end{figure}}
% Three panels in a row (small steps).
\newcommand{\wbthree}[5]{%
  \begin{figure}[htbp]\centering
    \begin{tabular}{@{}c@{\hskip3pt}c@{\hskip3pt}c@{}}
      \fig[0.30]{#1} &
      \fig[0.30]{#2} &
      \fig[0.30]{#3}
    \end{tabular}
    \caption{#4}\label{#5}
  \end{figure}}
% Four panels in a 2x2 grid.
\newcommand{\wbfour}[6]{%
  \begin{figure}[htbp]\centering
    \begin{tabular}{@{}c@{\hskip4pt}c@{}}
      \fig[0.44]{#1} &
      \fig[0.44]{#2} \\[2pt]
      \fig[0.44]{#3} &
      \fig[0.44]{#4}
    \end{tabular}
    \caption{#5}\label{#6}
  \end{figure}}
% A wide figure (benchmark grid etc.)
\newcommand{\wbwide}[3]{%
  \begin{figure}[htbp]\centering
    \includegraphics[width=\textwidth]{figures/#1}
    \caption{#2}\label{#3}
  \end{figure}}
% Pointer from the prose back to the live deck.
\newcommand{\deckslides}[1]{{\sffamily\footnotesize\textbf{Deck:} #1}}
% Inline code / command names.
\newcommand{\code}[1]{\texttt{\small #1}}
% A small "which files back this up" pointer.
\newcommand{\sources}[1]{{\footnotesize\textit{Sources:} #1}}
% Figure-panel label used inside captions.
\newcommand{\panel}[1]{\textbf{(#1)}}

\begin{document}

\markboth{Garcia-Dias}{Companion workbook: chemical tagging}

\title{Unsupervised Learning: High-dimensional embeddings (t-SNE, UMAP,
EVoC) and clustering algorithms applied to chemical tagging of star clusters}

\author{Rafael Garcia-Dias\affil{King's College London;
\texttt{rafaelagd@gmail.com}}}

\begin{abstract} Stars born in the same cluster share a chemical fingerprint,
and that fingerprint outlives the cluster: the stars disperse, their element
ratios do not. Recovering dispersed siblings from abundances alone
(\emph{chemical tagging}) is a clustering problem in a noisy, high-dimensional
space, and this workbook is the long-form companion to the IAA-SO 2026 lecture
on it. It builds the toolbox one algorithm at a time, each chosen to repair the
failure mode of the previous one: K-means, the $k$-nearest-neighbour primitive,
DBSCAN, HDBSCAN$^*$, PLSCAN, t-SNE, UMAP and EVoC, with the mathematics written
out, the hyperparameters interpreted, and the assumptions of each method stated
rather than hidden. The second half applies the toolbox to 25 open and globular
clusters in APOGEE DR19, Gaia DR3 and GALAH DR4, scoring every method against
kinematic membership, and reports the result the project actually found:
abundances narrow the candidate list, kinematics decide, and a masked spectral
autoencoder trained with no abundance labels at all separates clusters better
than the ASPCAP abundances do. Every number is reproducible from the
accompanying repository, and the failures (including one of our own, withdrawn)
are reported next to the successes.
\end{abstract}

\begin{keywords}
chemical tagging, clustering, dimensionality reduction, t-SNE, UMAP, EVoC,
stellar abundances, APOGEE, Gaia, unsupervised learning
\end{keywords}

\maketitle

\tableofcontents

% 01 · Introduction: the science question and the plan of the workbook
\section{INTRODUCTION}
\label{sec:intro}

\deckslides{slides 1--3 (the science question, the data, and the clustering
problem) and slide 51 for the reference list.}

Every star that exists today was born in a cloud of gas and dust, and most of
them were born in the company of thousands of siblings. A cluster forms from
one well-mixed reservoir, so its members inherit almost the same chemical
composition: element ratios that agree to better than 0.1 dex, far below the
spread between one cluster and the next. Then the cluster dissolves. Tidal
fields, encounters with giant molecular clouds and the shear of the Galactic
disc peel the members away one by one, and after a few hundred million years
(billions, for the most massive clusters) the stars are scattered across
kiloparsecs of the disc, moving on orbits that no longer have anything in
common. Their chemistry, however, is frozen. A star keeps the composition it
was born with, because nothing in a stellar interior efficiently mixes the
surface layers. \emph{Chemical tagging} is the attempt to use that frozen
fingerprint to put dispersed siblings back together: to look at a field of
stars, all of them currently unrelated, and work out which of them were born in
the same lost cluster \citep{Freeman:02, BlandHawthorn:16}.

This workbook is the long-form companion to the IAA-SO 2026 lecture on that
problem. The lecture is a tour of the unsupervised-learning toolbox with one
astronomical question attached to it; the workbook goes deeper on both halves:
the mathematics of each algorithm, and the measurements we made with it on 25
real open and globular clusters.

\begin{marginnote}[]
\entry{Chemical tagging}{Using the element abundances of stars to recover
which of them were born together, long after the birth cluster has dispersed.}
\entry{C-space}{The abundance space in which tagging happens: here 16
elements measured from a spectrum, one point per star.}
\end{marginnote}

\subsection{A cluster dissolves; its chemistry does not}

The premise is a physical fact about the interstellar medium: mixing is fast
and efficient, so a cluster of $10^3$--$10^5$ stars forms from gas that was
chemically homogeneous to a fraction of a dex. Those stars also formed in one
burst, so they share an age. What they do \emph{not} share, once the cluster
is gone, is position or velocity. Galactic archaeology therefore faces a
reconstruction problem: the cluster is a fossil, and chemistry is the only
part of it that is still legible.

Two things make the reconstruction hard, and both are properties of the data
rather than of the algorithms.

\begin{enumerate}
\item \textbf{Precision.} The abundance differences between clusters of similar
age and birth radius are small. At $0.05$--$0.1$ dex per element (roughly the
precision of a modern survey pipeline) the chemical differences between two
coeval solar-metallicity open clusters can be smaller than the measurement
noise for the elements in question \citep{GarciaDias:19, Casamiquela:21}.
Measured differentially (same instrument, same lines, cluster members only) the
internal coherence improves to $0.03$ dex, and the cluster-to-cluster
differences do not improve with it: the gas the clusters formed from was
already mixed to a comparable level, $0.02$--$0.03$ dex of scatter in nearby
spiral discs, correlated over scales below $600$\,pc \citep{Kreckel:20}. The
same efficient mixing that makes one cluster chemically coherent is what limits
how different two clusters can be.
\item \textbf{Contamination.} The field is crowded. In the sample used here,
  1\,002 stars belong to the 25 target clusters and 357\,056 do not: roughly
  one member in every 300 stars. Any search that returns a group has to be
  judged on both how many real members it found (recall) and how many of the
  stars in the group are real members (precision).
\end{enumerate}

\subsection{The task is clustering, not classification}

Chemical tagging is a \emph{clustering} problem: we are not handing labels to
a supervised model, we are asking an unsupervised algorithm to tell us how
many groups the data contains and which star belongs to which. That
distinction has consequences which reappear throughout the workbook.

\begin{issues}[WHY IT MATTERS THAT NOTHING IS LABELLED]
\begin{enumerate}
\item There is no ground truth inside the algorithm. A supervised classifier
  is scored on held-out labels; a clustering algorithm cannot be, so a
  clustering result is only as good as the validation you bring to it from
  outside (\textbf{Section~\ref{sec:validation}}).
\item Hyperparameters must be chosen without labels, which is why the method
  chapters spend as much time on \emph{what the knob means} as on the
  mathematics that uses it.
\item Feature choice is part of the model. In this problem the same algorithm
run on abundances, on a spectral latent, or on kinematics answers different
questions, and concatenating two such spaces does not combine their strengths,
it dilutes the smaller one (\textbf{Section~\ref{sec:benchmark}}).
\end{enumerate}
\end{issues}

\subsection{Two published verdicts, and the one this workbook adds}

The literature is not unanimous, and the disagreement is instructive because it
is not about methods. It is about tasks and precision.

\citet{Kos:17} ran t-SNE on GALAH abundances around the Pleiades and reported
success: seven of nine clusters recovered, plus two new Pleiades members six
degrees from the cluster centre. Their test case was a young cluster
($\sim\!100$ Myr), their verification was kinematic, and their grouping step
was visual inspection of the embedding.

\citet{GarciaDias:19} asked the same question with APOGEE abundances over
23 clusters and returned a sceptical answer: some clusters are chemically
indistinguishable from each other, and abundances alone are enough to
identify \emph{families} of clusters rather than individual ones. The
reconciliation is not that one study is wrong. The two tagged different
objects, at different ages, with different abundance precision, and asked
different things of the answer.

A third result belongs on the table too, because it cuts against the
pessimistic reading. \citet{Hogg:16} clustered 15 \emph{Cannon} abundances for
$\sim\!10^5$ APOGEE stars with K-means, using no positional information at all,
and showed that abundance-space overdensities really are phase-space clusters.
Their abundances reach $\sim\!0.04$ dex, comfortably inside the scatter of the
pipeline used here. The scope of any negative result in this field is therefore
\emph{these abundances, these clusters}, never ``abundances cannot tag''.

The current state of the art sits in between. \citet{Spina:25} trained a graph
attention autoencoder on chemistry, orbital similarity and age for
$\sim\!47\,000$ APOGEE thin-disc stars, obtained 282 groups and recovered five
of six open clusters. It is the natural competitor to this work and we have
not benchmarked against it; \textbf{Section~\ref{sec:spectral}} says what we
expect a like-for-like comparison to change.

\begin{marginnote}[] \entry{Strong tagging}{Recovering a specific lost birth
cluster from chemistry alone: the hard version, and the one this workbook
benchmarks.} \entry{Weak tagging}{Grouping stars into chemical families.
Populations that share a formation epoch rather than a common birth site. The
realistic outcome at current precision.}
\end{marginnote}

\subsection{What this workbook measures}

The second half is a benchmark we ran ourselves. We take 25 clusters (18 open,
7 globular) with kinematic membership from Gaia DR3, cut the APOGEE DR19 sample
around them, and ask three families of method (t-SNE, UMAP and EVoC) to recover
members from 16 abundances, from a 256-dimensional latent of the full spectrum,
or from kinematics, all scored against the same kinematic truth. On top of that
we train a masked spectral autoencoder that never sees an abundance label and
ask whether it tags better than the abundances themselves.

Four findings organise the results, each developed in its own section:

\begin{enumerate}
\item \textbf{Abundances narrow; kinematics decide.} Abundance-space
  clustering separates the 25 clusters only weakly (homogeneity $0.56$ for the
  best projector) where kinematics reach $0.94$, and adding abundances to
  kinematics never helps.
\item \textbf{Full-sky collapses; regions work.} Embedding $163\,000$ stars at
  once buries every cluster in field stars. The published t-SNE tagging
  operated on $30$--$45^\circ$ regions, and reproduces only there.
\item \textbf{The full spectrum beats the element ratios.} A self-supervised
latent of 8\,575 spectral pixels separates the same stars better than the 16
ASPCAP abundances, and the gap survives removing the globular cluster that
dominates the sample. A linear compression (PCA) sits in between, which is the
honest way to state the result.
\item \textbf{Chemical tagging is a two-survey problem.} APOGEE sees the
  giants that pin the distance; GALAH sees the main sequence that pins the
  age. Neither survey alone recovers both cluster parameters.
\end{enumerate}

\subsection{How to read this workbook}

Deck and workbook are complements. The deck carries the animations (K-means
converging, DBSCAN growing, the density sweep, the t-SNE descent, the UMAP
layout, the PLSCAN persistence barcode) and the workbook carries the equations,
the hyperparameter reasoning and the measured numbers. Where a deck figure
moves, the workbook prints a filmstrip of key frames and says so in the
caption.

\begin{figure}[htbp]\centering
  \begin{tabular}{@{}c@{\hskip18pt}c@{}}
    \fig[0.28]{qr_deck.png} &
    \fig[0.28]{qr_repo.png}
  \end{tabular}
  \caption{Left: the live deck, where the animated versions of the figures in
  this workbook run. Right: the repository that produced every number quoted
  in \textbf{Sections~\ref{sec:benchmark}} and \ref{sec:spectral}. Both are
  unlisted but openly reachable.}
  \label{fig:qrcodes}
\end{figure}

The chapters follow the order of the lecture. \S\ref{sec:data} is the data;
\S\ref{sec:fundamentals}--\S\ref{sec:evoc} are the toolbox, each method
motivated by the failure of the previous one: K-means
(\S\ref{sec:kmeans}), the $k$-nearest-neighbour primitive
(\S\ref{sec:knn}), DBSCAN (\S\ref{sec:dbscan}),
HDBSCAN$^*$ (\S\ref{sec:hdbscan}), PLSCAN (\S\ref{sec:plscan}), validation
(\S\ref{sec:validation}), t-SNE (\S\ref{sec:tsne}), UMAP
(\S\ref{sec:umap}) and EVoC (\S\ref{sec:evoc}).
\S\ref{sec:benchmark}--\S\ref{sec:physics} are the measurements, and
\S\ref{sec:assignment} is the assignment.

\begin{issues}[EXERCISES]
\begin{enumerate}
\item A cluster of $10^4$ stars dissolves over a few hundred Myr; a globular
  cluster of $10^6$ stars survives for more than 10 Gyr. Explain from that why
  chemical tagging should be \emph{easier} for old metal-poor clusters than
  for young solar-metallicity ones, then check your answer against the
  per-cluster results in \textbf{Section~\ref{sec:benchmark}}.
\item The field-to-member ratio here is about 300:1. If a search returns 100
  candidates for a cluster with 20 true members, what are its precision and
  recall? Which would you rather sacrifice, and what does that choice cost
  downstream?
\item Read \citet{Hogg:16} and \citet{Casamiquela:21} side by side: one is
  optimistic about strong chemical tagging and one is not. Write two sentences
  each on which \emph{measurement} they disagree about.
\end{enumerate}
\end{issues}

% 02 · The data — surveys, abundance space, kinematic truth, the cluster set
\section{THE DATA}
\label{sec:data}

\deckslides{slides 2 (the data) and 21 (the scale problem).}

Three surveys supply everything this workbook measures. APOGEE provides the
chemistry, Gaia the astrometry and radial velocities, and GALAH the southern
main sequence that appears in \S\ref{sec:physics}. The samples are built from
public data releases, and the whole pipeline --- download, quality cuts,
membership, matrices --- is reproducible with three commands
(\S\ref{sec:howtorun}).

\subsection{APOGEE: abundances for giants}

The Apache Point Observatory Galactic Evolution Experiment
\citep{Majewski:17} is a near-infrared
($H$-band, $1.51$--$1.70$ $\mu$m) spectroscopic survey of $\sim\!10^5$ stars
per data release, obtained with the 2.5-m Sloan Foundation Telescope in the
north and the du Pont Telescope in the south. The infrared window matters for
this project: it sees through the dust that hides the inner disc, and it
targets cool giants, which are bright enough to be observed across the whole
Galaxy.

Until recently the abundances came from the ASPCAP pipeline. The SDSS-V
release used here --- DR19 \texttt{astraAllStarASPCAP} \citep{Almeida:23},
$1.17$ GB, $1\,095\,480$ targets --- is produced by its successor, the
\emph{Astra} framework. The transfer is not cosmetic; two schema changes are
worth knowing about because they bite silently rather than loudly:

\begin{itemize}
\item abundances are published as \emph{$[$X/H$]$}, not $[$X/Fe$]$ under a
  different name, so element ratios must be formed explicitly as
  $[$X/Fe$] = [$X/H$] - [$Fe/H$]$;
\item the quality flags were renamed --- \texttt{ASPCAPFLAG} became
  \texttt{flag\_bad}, \texttt{STARFLAG} became \texttt{spectrum\_flags}.
\end{itemize}

DR19 also reanalyses rather than replaces DR17: 717\,689 of its rows carry
\texttt{release='dr17'}, all reprocessed by a single version of the pipeline.
That detail becomes important in \textbf{Section~\ref{sec:spectral}}, where an
earlier version of our benchmark was nearly derailed by mixing data
\emph{products} from the two.

\begin{marginnote}[]
\entry{ASPCAP / Astra}{The abundance pipelines: ASPCAP up to SDSS-IV DR17,
Astra from SDSS-V DR19 onwards. Same philosophy --- fit synthetic spectra
to the observation --- different catalogues and column names.}
\entry{Abundance precision}{The per-element statistical error on $[$X/H$]$,
typically $0.02$--$0.1$ dex for APOGEE giants, depending on the element and
the signal-to-noise ratio. It sets the ceiling on what chemistry can tag.}
\end{marginnote}

\subsection{Gaia and GALAH: the referee and the main sequence}

Gaia DR3 \citep{Gaia:23} supplies parallaxes, proper motions and the
cross-matched photometry for every star in the sample, all of it carried
inside the APOGEE catalogue for convenience. These are the quantities that
define membership (\S\ref{sec:kinematics}); we never feed them to a
tagging method, and \textbf{Section~\ref{sec:validation}} explains why that
asymmetry is the whole point.

GALAH DR4 \citep{DeSilva:15, Buder:21} covers declination $\lesssim +25^\circ$
with a high-resolution optical spectrograph, which means it sees the
\emph{main sequence} where APOGEE sees giants. It enters only in
\S\ref{sec:physics}, where the two surveys are combined to recover a cluster
age and distance at the same time.

\subsection{The abundance space}

Everything is standardised into one matrix: 16 numbers per star, one column
per element, built from the survey catalogue by a fixed sequence of steps.
Each step is a modelling choice, not housekeeping, and each has a failure mode
we have hit.

\begin{enumerate}
\item \textbf{Quality cuts.} Signal-to-noise ratio $\geq 100$, ASPCAP flag
  zero, stellar flag zero. Cuts this aggressive remove the noisy tail at the
  cost of sample size; loosening them changes the answer
  (\S\ref{sec:benchmark}, the levers).
\item \textbf{The element set.} C, N, O, Na, Mg, Al, Si, S, K, Ca, Ti, V,
  Cr, Mn, Ni --- the $\alpha$ and iron-peak elements --- plus $[$Fe/H$]$ as
  the sixteenth coordinate. This is the set used by the reference papers, so
  results are comparable with them.
\item \textbf{Missing values.} Weak lines go undetected in metal-poor stars
  and their abundances come back as NaN. Dropping every star with a NaN
  deletes the metal-poor globulars \emph{entirely} --- M~15 and M~92 vanish
  from the sample, and with them the cleanest tagging cases in the whole
  project. We therefore keep stars with at least 8 of 16 finite elements and
  fill the remaining gaps with the column median. This single choice is the
  difference between a sample that contains the interesting objects and one
  that does not.
\item \textbf{Standardisation.} Each element is shifted to zero median and
  scaled to unit standard deviation, so that no element dominates the distance
  metric by virtue of its units. (An early bug in this project applied
  inverse-variance weights \emph{before} this step, which silently cancelled
  them; see \S\ref{sec:benchmark}.)
\item \textbf{Row normalisation.} Each star's 16-vector is scaled to unit
  $L_2$ norm. The reason is metric compatibility: after this step, Euclidean
  distance is a monotone function of cosine distance, so t-SNE and UMAP (which
  use Euclidean) and EVoC (whose native metric is cosine) all see the same
  geometry, and the comparison between them is fair. It is also the single
  most effective precision lever in the pipeline: without it, density-based
  clustering merges the whole field into one blob.
\end{enumerate}

\begin{figure}[htbp]\centering
  \fig[0.62]{cspace.png}
  \caption{The shape of the problem, in two of the sixteen dimensions (from
  the deck). Field stars in grey; the members of one cluster in red. The
  cluster is a tight, slightly elongated clump embedded in a cloud several
  orders of magnitude larger --- a small overdensity in a big, structured
  background. Everything in
  \S\ref{sec:fundamentals}--\S\ref{sec:evoc} is a strategy for finding
  objects like this one.}
  \label{fig:cspace}
\end{figure}

\subsection{Kinematic membership is the referee, not a feature}
\label{sec:kinematics}

To score anything we need labels, and this is where chemical tagging is
unusually lucky: a cluster is also a kinematic object. Its members currently
share a position on the sky, a parallax, a proper motion and a radial
velocity, because they were born at the same place with almost the same
velocity. We use those four observables, $\sigma$-clipped around the cluster
centre, to label each star \emph{member} or \emph{field}.

The circularity risk is obvious: if kinematics are used to build the labels,
then a method that receives kinematics as an input is being scored on its own
input. We therefore treat kinematics as a \emph{ceiling}, not a competing
signal --- the number that says how well \emph{any} chemistry-based method
could possibly do on this sample. Published cluster catalogues provide an
independent cross-check: scoring against Simbad memberships gives numbers
within a few percent of the kinematic ones, and both are reported in the
benchmark \citep{Dias:02, Harris:96}.

\subsection{The cluster set}
\label{sec:clusters}

The target list is the 23 clusters of \citet{GarciaDias:19} plus the Pleiades
of \citet{Kos:17} --- and, for the two-survey work, the southern open clusters
NGC~2243 and Collinder~261. Twenty-five objects in total: seven globular and
eighteen open. They span two orders of magnitude in age, from the Pleiades at
$\sim\!100$ Myr to NGC~188 at $7.6$ Gyr, and two orders of magnitude in mass.
The counts below come from the full DR19 run (357\,056 field stars). They are
\emph{row counts}, not star counts: the sample contains repeat entries for the
same star, a defect quantified in \S\ref{sec:validation} and one we flag here
rather than at the end because it affects how every macro average in this
workbook should be read.

\begin{table}[htbp]
\tabcolsep5pt
\caption{The 25 target clusters and the number of kinematic members each
contributes in the DR19 sample. Globulars are marked with a $\dagger$. The
asymmetry in the counts is real and consequential: M~67 alone supplies
$23\%$ of all members, while three clusters supply fewer than ten.}
\label{tab:clusters}
\begin{center}
\begin{tabular}{@{}lrlr@{}}
\hline
Cluster & $n_{\rm mem}$ & Cluster & $n_{\rm mem}$\\
\hline
Berkeley 17 & 7 & M 3$^{\dagger}$ & 154\\
Berkeley 66 & 12 & M 5$^{\dagger}$ & 67\\
Berkeley 71 & 14 & M 13$^{\dagger}$ & 34\\
Collinder 261 & 22 & M 15$^{\dagger}$ & 31\\
IC 166 & 35 & M 71$^{\dagger}$ & 39\\
King 5 & 12 & M 92$^{\dagger}$ & 25\\
King 7 & 8 & NGC 1245 & 21\\
M 67 & 230 & NGC 1798 & 13\\
Pleiades & 23 & NGC 188 & 25\\
& & NGC 2158 & 6\\
& & NGC 2243 & 36\\
& & NGC 2420 & 19\\
& & NGC 6791 & 45\\
& & NGC 6819 & 62\\
& & NGC 7789 & 47\\
\hline
\multicolumn{4}{@{}l}{M 107$^{\dagger}$ contributes 15 members.\quad
Total: 1\,002 members in 25 clusters.}\\
\hline
\end{tabular}
\end{center}
\begin{tabnote}
Counts from the full-sky DR19 run (\texttt{cluster run}), field star cap
disabled, on the member frame the DR19 baseline scores --- every member row
with at least nine finite abundances out of sixteen. Reproduce with
\texttt{python scripts/population\_counts.py --baseline-min-finite 9};
the counts are identical for any floor between nine and fifteen. Member counts
differ slightly between runs because the clustering must first find a group
before its overlap with the truth can be measured; these are the truth-side
counts.
\end{tabnote}
\end{table}

\begin{figure}[htbp]\centering
  \begin{tabular}{@{}c@{\hskip8pt}c@{}}
    \fig[0.42]{sky_m67.png} &
    \fig[0.42]{sky_m3.png}
  \end{tabular}
  \caption{Two of the targets, in the optical (DSS2 colour composites centred
  on the cluster): M~67 (left), an old solar-metallicity open cluster, and
  M~3 (right), a metal-poor globular. The two objects illustrate the range of
  the sample and, in the results, the range of difficulty: the globular is
  chemically distinct from the field around it, the open cluster is not.}
  \label{fig:skysamples}
\end{figure}

\subsection{How to reproduce the sample}
\label{sec:howtorun}

The pipeline is one package with a single configuration file, so every number
in this workbook can be regenerated from a fresh checkout. It is written on
\texttt{numpy} \citep{Harris:20}, \texttt{pandas} \citep{McKinney:10},
\texttt{scikit-learn} \citep{Pedregosa:11}, \texttt{astropy}
\citep{Astropy:22} and \texttt{matplotlib} \citep{Hunter:07}, with the
clusterers taken from their reference implementations. The commands below
are all you need to rebuild the sample and the benchmark.

\begin{summary}[HOW WE RAN IT]
\begin{quote}\small
\begin{verbatim}
git clone github.com/garciadias/iaa-advanced-neural-networks-2026-draft
cd iaa-advanced-neural-networks-2026-draft
uv sync                       # Python 3.13 environment
uv run cluster download       # 1.17 GB, SDSS-V DR19 Astra ASPCAP
uv run cluster run --fast     # ~1 min on a laptop, capped field
CLUSTER_FAST=0 uv run cluster run   # full all-sky run, 10-20 min
\end{verbatim}
\end{quote}
Every flag lives in \texttt{src/cluster/config.py} and every flag also has an
environment override with the \texttt{CLUSTER\_} prefix, which is how the
variants in \S\ref{sec:benchmark} were produced. The Docker path
(\texttt{docs/docker.md}) exists for machines without \texttt{uv}.
\end{summary}

\begin{table}[htbp]
\centering
\small
\tabcolsep4pt
\begin{tabular}{@{}ll@{}}
\toprule
setting & value used throughout \\
\midrule
\texttt{SNR\_MIN} & $100$ \\
dwarf cut (\texttt{DWARF\_ONLY}) & off \\
elements & 16: \texttt{[C/Fe]} \dots\ \texttt{[Ni/Fe]} (15 ratios) $+$ \texttt{[Fe/H]} \\
imputation & median, requiring $\geq\!8$ finite elements \\
standardisation & on (zero median, unit s.d.) \\
row normalisation & on \\
element weights & off \\
\texttt{HDBSCAN.min\_cluster\_size} & $5$ \\
\texttt{TSNE.perplexity} / \texttt{UMAP.n\_neighbors} & $30$ / $15$ \\
\texttt{random\_state} & $42$ \\
\bottomrule
\end{tabular}
\caption{The pipeline's settings for every number in this workbook, unless a
variant is named. Overrides are environment variables with the
\texttt{CLUSTER\_} prefix --- \texttt{CLUSTER\_NORMALIZE\_ROWS=0} is the one
that most changes the answer.}
\label{tab:config}
\end{table}

\begin{issues}[EXERCISES]
\begin{enumerate}
\item Take the 16-element vector for one star in the sample. Standardise it
  and then row-normalise it. Show algebraically that after row normalisation
  the Euclidean distance between two stars is a strictly monotone function of
  the cosine distance between them. Why does EVoC care?
\item \texttt{MIN\_FINITE\_ELEMENTS} defaults to 8 of 16. Set it to 16 (strict
  complete-case) and count how many of the 25 clusters survive with at least
  five members. Which clusters are lost, and what does that do to any
  conclusion about the surviving ones?
\item M~67 contributes 230 of 1\,002 members. Recompute a macro-average
  metric over clusters with and without M~67 in the pool. How much of the
  published summary depends on one cluster?
\end{enumerate}
\end{issues}

% 03 · Clustering fundamentals — what a cluster is, and what breaks
\section{CLUSTERING FUNDAMENTALS}
\label{sec:fundamentals}

\deckslides{slides 3 (the clustering problem), 18 (no free lunch) and 19--20 (validation).}

Before any algorithm, two questions deserve a straight answer: what do we
mean by ``a cluster'', and what does a distance in 16 dimensions actually
measure? The answers are not academic. Most of the failures in
\textbf{Section~\ref{sec:benchmark}} are failures of an assumption about one
of those two things, not failures of an implementation.

\subsection{Four definitions of ``cluster''}

A clustering algorithm does not discover the clusters that are in the data;
it discovers the clusters that its own definition of ``cluster'' implies. The
definitions in play in this workbook are:

\begin{itemize}
\item \textbf{Centroid.} A cluster is a set of points closer to one centre
  than to any other. This is K-means (\S\ref{sec:kmeans}), and it implies
  convex, roughly equal-sized groups.
\item \textbf{Density.} A cluster is a connected region where the local point
  density exceeds a threshold. This is DBSCAN and its descendants
  (\S\ref{sec:dbscan}--\S\ref{sec:plscan}), and it implies that clusters and
  noise are distinguishable by density alone.
\item \textbf{Connectivity.} A cluster is a connected component of a graph
  built from nearest neighbours, with edges weighted by similarity. This is
  the graph view that t-SNE, UMAP and EVoC share
  (\S\ref{sec:tsne}--\S\ref{sec:evoc}).
\item \textbf{Statistical.} A cluster is a component of a mixture model whose
  parameters are inferred by maximum likelihood (the EM algorithm,
  \citealp{Dempster:77}). The cluster boundaries are then soft: a star has a
  probability of belonging to each group.
\end{itemize}

The reference chapter behind the lecture's framing --- \citet{GarciaDias:20},
a methods chapter written for a psychiatry audience --- makes the same point
from the other direction: clustering is worth attempting when the joint
distribution of the variables is \emph{multipeaked}, and the first decision is
whether the data are multimodal at all. If a distribution is unimodal, every
algorithm will return groups anyway, and every one of them will be an
artefact.

\begin{marginnote}[]
\entry{Multimodality}{A distribution with more than one peak. Clustering is
only meaningful if the peaks are physical --- here, if stars of different
birth sites occupy different regions of C-space.}
\entry{No free lunch}{No algorithm is best on every problem: each encodes
assumptions about shape, scale, density and metric. Choosing a method is
choosing which assumptions are safe for your data.}
\end{marginnote}

\subsection{Distance in C-space}

Every algorithm in this workbook reduces to a similarity judgement, and all of
them use either Euclidean or cosine distance. Two transformations in
\textbf{Section~\ref{sec:data}} make those two compatible, and one
transformation is a genuine modelling choice:

\[
d_{\rm Euclid}(\mathbf{x},\mathbf{y}) = \left\| \mathbf{x} - \mathbf{y}
\right\|_2, \qquad
d_{\rm cos}(\mathbf{x},\mathbf{y}) = 1 - \frac{\mathbf{x}\cdot\mathbf{y}}
{\|\mathbf{x}\|_2\,\|\mathbf{y}\|_2}.
\]

After per-element standardisation and per-star $L_2$ normalisation, all stars
lie on the unit sphere in 16 dimensions and $d_{\rm Euclid}$ is a monotone
function of $d_{\rm cos}$ --- so t-SNE and UMAP (Euclidean internally) and
EVoC (cosine) are looking at the same geometry, and comparing them is fair.
The Euclidean alternative without normalisation is not wrong, but it inherits
a scale from the standardisation: elements with a long tail of outliers become
the elements that decide every distance.

High dimension does the rest. Two effects matter here.

\begin{enumerate}
\item \textbf{Concentration.} In $d$ dimensions, the ratio between the
  furthest and nearest neighbour distance shrinks as $d$ grows: distances
  \emph{concentrate} around a common value. A density threshold that worked in
  2-D has to be re-tuned, and the contrast between ``near'' and ``far''
  degrades. This is the everyday experience of the curse of dimensionality in
  clustering: the kNN graph of 16-D abundances has edges that are only
  slightly shorter than the edges you would get by chance.
\item \textbf{Dilution by concatenation.} Adding dimensions that carry no
  cluster information can only reduce the relative contribution of the
  dimensions that do. The consequence in this project is concrete. On the same
  982 member stars, in the same pipeline
  (\textbf{Section~\ref{sec:benchmark}}), the cluster-only homogeneity is
  $0.56$ from 16 abundances, $0.94$ from 4 kinematic coordinates, and
  $0.75$ from the 20-dimensional concatenation of the two. The blend is
  \emph{worse} than the kinematics alone: the chemistry did not add
  information, it added dimensions.
\end{enumerate}

That last result deserves emphasis because the intuitive fix --- ``if both
signals are useful, feed both'' --- is the one most people try first:

\begin{issues}[THE TRAP OF CONCATENATION]
Fusing two signals of very different dimensionality into one Euclidean space
is not a fusion of information; it is a weighted sum in which the weight of
each channel is set by how many coordinates it happens to contribute. A
256-dimensional spectral latent next to 4 kinematic coordinates gives the
spectrum 98\% of the vote, whatever the physics says. The two ways out are to
reduce the high-dimensional channel before concatenating it, or to combine the
two at the level of the \emph{decisions} rather than the features --- the
two-stage pipeline of \S\ref{sec:benchmark}, which uses kinematics for
recall and a spectral distance for precision.
\end{issues}

\subsection{Four tasks that are easily confused}

The literature --- and our own first drafts --- mixes up four different
questions that all get called ``clustering''. Keeping them separate is the
single most useful discipline in this project.

\begin{enumerate}
\item \textbf{Membership determination.} For a given cluster, label each star
  in a field \emph{member} or \emph{field}. This is the operationally useful
  task, and it is what a survey pipeline actually needs.
\item \textbf{Discovery.} Blind search: find groups without knowing in advance
  where to look or what they should look like. This is what chemical tagging
  means in the ``strong'' sense of \S\ref{sec:intro}, and it is much harder:
  the algorithm must invent the groups as well as find their members.
\item \textbf{Refinement.} Clean an existing candidate list --- remove field
  contaminants from a known cluster's sample. This is the task our two-stage
  pipeline performs, and it is measurably easier than discovery.
\item \textbf{Cluster-only separation.} Take the members of several clusters,
  with no field stars present, and ask whether the clusters are
  distinguishable from one another. This is the experiment of
  \citet{GarciaDias:19}, recreated in \S\ref{sec:benchmark} as the
  \emph{baseline}. Its score is an upper bound on what any field-level search
  could do: if clusters cannot be told apart when nothing else is present,
  they cannot be recovered from a field either.
\end{enumerate}

\begin{table}[htbp]
\tabcolsep6pt
\caption{The assumptions each method in this workbook makes explicit. Read it
as a checklist: when a method fails on our data, one of these entries is
usually the reason, and the next chapter takes over because of it.}
\label{tab:assumptions}
\begin{center}
\begin{tabular}{@{}p{4.6pc}p{8.2pc}p{9.6pc}@{}}
\hline
Method & Assumes & Fails when\\
\hline
K-means & clusters are convex, similar in size, separable by a centre;
  $K$ known & shapes are elongated or nested; outliers pull centres\\
$k$NN & a metric is meaningful and local density is informative &
  dimensionality is high; the $k$ is chosen without looking at the data\\
DBSCAN & one density threshold separates cluster from noise &
  clusters and noise share a density; densities vary between clusters\\
HDBSCAN$^*$ & density varies, but a single minimum cluster size applies &
  clusters are smaller than the floor; the metric is misleading\\
PLSCAN & persistence, not size, identifies a cluster &
  clusters are genuinely short-lived in scale-space (rare)\\
t-SNE & local neighbourhoods are what matter &
  inter-cluster distances are being interpreted\\
UMAP & the manifold is locally uniform &
  global geometry is treated as metric information\\
EVoC & a cosine-style graph and one internal embedding serve all scales &
  the metric is wrong for the data\\
\hline
\end{tabular}
\end{center}
\end{table}

\subsection{Validation, in outline}

Because there are no labels, validation is a design decision, not a
calculation. The toolbox has two halves, developed in
\textbf{Section~\ref{sec:validation}}:

\begin{itemize}
\item \textbf{Internal} indices --- sum of squared errors, silhouette, the dip
  test --- measure how \emph{tidy} a partition is. They are computed from the
  data alone, they are cheap, and they cannot tell you whether the partition
  is real.
\item \textbf{External} indices compare the partition with something the
  algorithm never saw. In this project that something is kinematic membership;
  in general it is any independent measurement.
\end{itemize}

The dangerous middle is a number that looks external but is not: a metric that
uses the quantity you are trying to predict, or a baseline that has collapsed
to two groups and scores well for the wrong reason. Both appear in our own
results and are flagged there.

\begin{issues}[EXERCISES]
\begin{enumerate}
\item The concentration of distances can be demonstrated in five lines:
  sample $n=1000$ points uniformly in $\mathbb{R}^d$, compute for each point
  the ratio of the distance to its nearest and to its furthest neighbour, and
  plot the mean ratio against $d$ for $d = 1,\dots,50$. At what $d$ does the
  nearest neighbour stop being meaningfully nearer than the furthest one?
\item Standardise the 16-D matrix, then row-normalise it. A star with an
  unusually low $[$Fe/H$]$ and a star with an unusually high one now have the
  same norm. What did normalisation do to the information about overall
  metallicity, and why is that mostly a feature here rather than a bug?
\item Pick three clusters from \textbf{Table~\ref{tab:clusters}} --- one with
  more than 60 members, one with about 20, one with fewer than 10 --- and
  write down what ``$K$ is known'' would even mean for a field of 25 clusters
  with these counts. Which of the four tasks of \S 3.3 does a fixed $K$
  violate?
\end{enumerate}
\end{issues}

% 04 · K-means — the default tool, and its assumptions
\section{K-MEANS: THE DEFAULT TOOL}
\label{sec:kmeans}

\deckslides{slides 4--11 --- the K-means tour, ending on its limits.}

K-means is where every clustering analysis starts, and it is worth
understanding precisely, because its failure modes are the ones the rest of
the toolbox was invented to repair. The method was proposed by
\citet{MacQueen:67} and its modern formulation --- the one everyone implements
--- is due to \citet{Lloyd:82}.

\subsection{The objective}

Given $n$ points $\mathbf{x}_1,\dots,\mathbf{x}_n$ in $\mathbb{R}^d$ and a
fixed number of groups $K$, K-means looks for $K$ centres
$\boldsymbol{\mu}_1,\dots,\boldsymbol{\mu}_K$ and an assignment of every point
to one centre that minimises the \emph{sum of squared errors} (SSE), also
called the within-cluster sum of squares:

\begin{equation}
J = \sum_{k=1}^{K} \sum_{\mathbf{x} \in C_k}
\left\| \mathbf{x} - \boldsymbol{\mu}_k \right\|_2^2 ,
\qquad
\boldsymbol{\mu}_k = \frac{1}{|C_k|}\sum_{\mathbf{x} \in C_k} \mathbf{x}.
\label{eq:sse}
\end{equation}

The second half of the equation is not a separate model choice but a
consequence of the first: for a fixed assignment, the centre that minimises
the sum of squared distances to the members is the arithmetic mean. Two
properties follow immediately and both are assumptions in disguise.

\begin{itemize}
\item Because $J$ penalises distance \emph{squared}, the objective is
  dominated by the points furthest from their centre. A single outlier can
  drag a centre across the field. K-means is not robust; medians (K-medians)
  or medoids are the standard remedy when outliers are expected.
\item Because every point is assigned to its nearest centre, the partition is
  a \emph{Voronoi tessellation}: cells are convex and their boundaries are
  hyperplanes. A cluster with a crescent or dumbbell shape cannot be
  represented, however obvious it looks to the eye
  (\textbf{Figure~\ref{fig:kmeansfail}}).
\end{itemize}

\begin{marginnote}[]
\entry{SSE / inertia}{The sum of squared distances from each point to its
assigned centre: the quantity K-means minimises. Also called inertia in
scikit-learn and within-cluster sum of squares (WCSS).}
\entry{Voronoi cell}{The region of space closer to one centre than to any
other. K-means partitions the plane into Voronoi cells by construction.}
\end{marginnote}

\subsection{The algorithm, and why it always stops}

Lloyd's algorithm is a two-move loop. \emph{Assign} every point to its nearest
centre; \emph{update} every centre to the mean of its members; repeat. Each
move can only reduce $J$: the assignment step is optimal for fixed centres,
and the update step is optimal for fixed assignment. $J$ is bounded below by
zero and takes finitely many values (there are finitely many partitions of
$n$ points into $K$ groups), so the loop must terminate --- in practice after
a few dozen iterations. It terminates at a \emph{local} minimum, and the
distinction matters: K-means is a coordinate-descent algorithm on a
non-convex objective, so different starting centres end in different answers.

\begin{figure}[htbp]\centering
  \fig[0.94]{kmeans_film.png}
  \caption{The K-means loop on a synthetic 2-D dataset, sampled at four
  points in the animation: initial centres, first assignment, first update,
  converged. The workbook prints the key frames; the deck plays the full
  sequence (\textbf{Figure~\ref{fig:qrcodes}}).}
  \label{fig:kmeansfilm}
\end{figure}

\begin{figure}[htbp]\centering
  \begin{tabular}{@{}c@{\hskip6pt}c@{}}
    \fig[0.44]{kmeans_step1.png} &
    \fig[0.44]{kmeans_step2.png} \\[3pt]
    \fig[0.44]{kmeans_step3.png} &
    \fig[0.44]{kmeans_step4.png}
  \end{tabular}
  \caption{The same four states, in the deck's step figures. Stars are the
  centres, dots the data: \panel{a} all points unassigned with $K$ centres
  dropped at random; \panel{b} assignment to the nearest centre; \panel{c}
  centres moved to the mean of their members; \panel{d} the fixed point of the
  iteration --- no point changes hands, so the loop stops.}
  \label{fig:kmeanssteps}
\end{figure}

\subsection{Choosing $K$, and choosing where to start}

$K$ is an input, which is uncomfortable in a problem where the number of
clusters is exactly what we are trying to learn. Three standard answers, and a
warning:

\begin{enumerate}
\item \textbf{Elbow.} Plot $J$ against $K$; $J$ decreases monotonically
  (more centres always fit better), and the curve is expected to bend where
  extra groups stop buying real structure. It is a heuristic: on smooth data
  there is no elbow, and on clustered data with a broad size distribution the
  bend is shallow.
\item \textbf{Silhouette.} For each point, compare its mean distance to its own
  cluster with its mean distance to the nearest other cluster; average over
  points and pick the $K$ that maximises the score \citep{Rousseeuw:87}. It is
  a ratio of within- to between-cluster distance, so it also penalises
  non-convex clusters.
\item \textbf{Model selection.} Compare $K$ by held-out likelihood under a
  Gaussian mixture fitted by EM \citep{Dempster:77}. K-means is the hard
  assignment limit of that model with spherical, equal-variance components ---
  which is precisely the assumption list of \textbf{Table~\ref{tab:assumptions}}.
\end{enumerate}

The warning is that in this project we \emph{cheat}, deliberately. The
catalogue tells us there are 25 clusters, so the benchmarking arms are handed
the true number of groups. That makes the comparison between methods clean,
and it flatters all of them relative to a blind search --- which is why
\S\ref{sec:benchmark} reports field retrieval (a discovery-like task) and
cluster-only separation (a $K$-known task) separately.

Initialisation has the same status: a modelling decision that changes the
answer. The classical approach is to run the loop from several random starts
and keep the best $J$; the modern default is \emph{k-means++}
\citep{Arthur:07}, which samples the first centre uniformly and each
subsequent centre with probability proportional to $D(\mathbf{x})^2$, the
squared distance to the nearest centre already chosen. The effect is to spread
the initial centres out, which typically reaches a better local optimum in
fewer restarts.

\begin{figure}[htbp]\centering
  \fig[0.58]{kmeans_init.png}
  \caption{Initialisation sensitivity (from the deck): the same data, two
  random starts, two different final partitions --- and one of them splits a
  real group. Nothing about the algorithm changed; only the seed did. In a
  clustering paper this is the first thing a referee will ask about, and it is
  why every number in this workbook is reported as a mean over seeds (seven,
  in the head-to-head of \S\ref{sec:spectral}).}
  \label{fig:kmeansinit}
\end{figure}

\subsection{What it assumes, and when it breaks}

\begin{issues}[THE ASSUMPTION LIST]
K-means assumes that clusters are (i) separable by a centre, (ii) roughly
convex, (iii) of comparable size and variance, (iv) all present in equal
proportion, and (v) that $K$ is known. It additionally assumes that the metric
is meaningful --- that the dimensions have been standardised, or that the
scales carry physical meaning. When it fails, it fails in a characteristic
way: it splits large or elongated groups and merges neighbouring small ones,
and it never returns ``noise''.
\end{issues}

\begin{figure}[htbp]\centering
  \fig[0.50]{kmeans_fail.png}
  \caption{K-means on two elongated, interleaved structures. The natural
  groups are obvious to the eye and unavailable to a partition by nearest
  centre: the objective can only be reduced by cutting across the arms. This
  is the failure that motivates every method in the rest of the workbook.}
  \label{fig:kmeansfail}
\end{figure}

\begin{figure}[htbp]\centering
  \fig[0.90]{kmeans_storyboard.png}
  \caption{The full K-means story in one grid, as used in the lecture: the
  loop, an initialisation that goes wrong, the failure on non-convex shapes,
  and the diagnostic that the objective can be lowered without the answer
  getting better.}
  \label{fig:kmeansstory}
\end{figure}

\subsection{Is K-means useful for chemical tagging?}

Yes --- with the right expectations, and in the literature it has been used
twice in ways that matter here.

\citet{Hogg:16} clustered $\sim\!10^5$ APOGEE stars in a 15-dimensional
abundance space with K-means, with no positional information, and recovered
groups that are genuinely phase-space structures. That result defines what
abundance-space overdensities can be at $\sim\!0.04$ dex precision. And
\citet{GarciaDias:18} applied K-means to $153\,847$ normalised APOGEE
spectra for spectral classification rather than for tagging, and the classes it
returns do track the stellar parameters --- dwarfs separate from giants, bulge
from halo. But the same paper's conclusion is the opposite of tidy: flux space
has no obvious groups, a discrete classification does not organise parameter
space neatly, and the solution moves with the initialisation. That is the same
warning as the row-order result of \S\ref{sec:validation}, in a different space.
K-means is useful when the groups are blobs; the spectra are not blobs, and the
abundance space of \S\ref{sec:benchmark} is only partly one.

In \textbf{Section~\ref{sec:benchmark}} the same algorithm is run on the same
member matrix as the other arms --- the 1\,002-row matrix of
\textbf{Table~\ref{tab:clusters}} in this chapter's summary box, and the
982-star deduplicated population of \textbf{Table~\ref{tab:honest}} and
\textbf{Table~\ref{tab:literature}} --- so that the comparison between ``the
classical tool'' and the newer ones is measured rather than asserted.

\begin{summary}[K-MEANS ON OUR MEMBER MATRIX]
\texttt{article/scripts/kmeans\_baseline.py} loads the DR19 sample, takes the
same member-only matrix the other benchmark arms use ($1\,002$ stars, 16
abundances, 25 clusters, $\geq\!5$ members), runs K-means with $n_{\rm
init}=10$ and reports the paper's merit functions alongside the SSE curve. The
measured result is worth stating plainly, because it is not the one the gospel
predicts: at $K=25$, K-means reaches \textbf{homogeneity 0.449}
(completeness 0.493, v-measure 0.470, accuracy 0.296) --- higher than t-SNE's
0.292 and EVoC's 0.409 on the same stars, and within reach of UMAP's 0.547.
Two caveats travel with it. At $K=2$ the same matrix gives homogeneity 0.228
with completeness 0.954: the ``families'' reading, in which almost everything
lands in one group. And raising $K$ to 40 lifts homogeneity to 0.505 while
completeness falls to 0.474 --- over-splitting buys homogeneity, which is why
the paper's metric pair must always be quoted together.
\end{summary}

\begin{issues}[EXERCISES]
\begin{enumerate}
\item Show that the assignment and update steps each cannot increase $J$, and
  conclude that the algorithm terminates. Then construct a one-dimensional
  example (three points, two centres) where the global optimum and Lloyd's
  fixed point differ, and draw the two partitions.
\item Implement k-means++ seeding in ten lines of NumPy and compare the final
  $J$ over 50 random datasets against a uniform-random initialisation. How
  much of the improvement comes from the seeding rule, and how much from the
  multiple restarts?
\item K-means with $K=2$ on our member matrix asks a different question from
  $K=25$: not ``which cluster is this star in'' but ``which family''. Run
  both and compare the homogeneity. Which of the two numbers would you quote
  to support strong chemical tagging, and why is that the wrong number?
\item Replace the Euclidean distance in the assignment step with cosine
  distance, keeping the mean update. What happens to the objective's
  guarantee of monotone decrease, and does the algorithm still terminate?
\end{enumerate}
\end{issues}

% 05 · KNN — the nearest-neighbour primitive
\section{THE kNN PRIMITIVE}
\label{sec:knn}

\deckslides{slide 12, with slide 23 for the core distance $\kappa(x)$.}

Before density-based clustering and before graph embeddings, there is one
operation that all of them depend on: find a point's $k$ nearest neighbours.
It is the oldest idea in the toolbox \citep{Cover:67}, it is simple enough to
be taken for granted, and in this project it turns out to be both a
clustering ingredient and a validation metric in its own right.

\subsection{Neighbours as a classifier}

The classical $k$-nearest-neighbour rule is a supervised method: to classify a
point, find its $k$ nearest labelled neighbours and take a vote. Its appeal
here is diagnostic rather than predictive. A $k$NN vote hides no assumptions
about the shape of a decision boundary, so an error rate measured with $k$NN
approximates the best you can do with a local rule in that space. When we ask
whether 16 abundances carry cluster information at all
(\S\ref{sec:benchmark}), a nearest-neighbour view answers the question
directly: do members share their neighbourhood with other members?

\begin{marginnote}[]
\entry{Core distance}{$\kappa_k(\mathbf{x})$, the distance from a point to its
$k$-th nearest neighbour. Large $\kappa$ means an isolated point; small
$\kappa$ means a dense neighbourhood.}
\entry{kNN graph}{The directed graph in which every point is joined to its $k$
nearest neighbours. Its nodes are stars, its edges similarity, and it is the
substrate that HDBSCAN$^*$, UMAP and EVoC all build on.}
\end{marginnote}

\subsection{The k-th neighbour distance is a density estimate}

Distances and densities are two descriptions of the same local geometry. The
volume of the ball containing exactly a fraction $k/n$ of the data is
$V \propto \kappa_k(\mathbf{x})^d$ in $d$ dimensions, so a local density
estimate follows from the measured distance alone:

\begin{equation}
\hat{\rho}(\mathbf{x}) \;\propto\; \frac{k}{n\,\kappa_k(\mathbf{x})^{d}}.
\end{equation}

Two consequences are used repeatedly in the chapters that follow. First,
$\kappa$ converts a distance into a density, which is the operation DBSCAN
performs with its global $\varepsilon$ and HDBSCAN$^*$ performs locally. Second,
the conversion involves $d$, the dimension: in 16 dimensions a factor of two
in density is a mere 4\% change in $\kappa$, which is why density-based
methods on abundance data are so sensitive to the metric and to the
normalisation chosen in \S\ref{sec:data}.

\begin{figure}[htbp]\centering
  \begin{tabular}{@{}c@{\hskip6pt}c@{}}
    \fig[0.42]{knn_step1.png} &
    \fig[0.42]{knn_step2.png} \\[3pt]
    \fig[0.42]{knn_step3.png} &
    \fig[0.42]{knn_step4.png}
  \end{tabular}
  \caption{The primitive, in four steps (deck figures): \panel{a} measure every
  pairwise distance from the query point (black); \panel{b} keep the $k$
  smallest --- the neighbourhood; \panel{c} the $k$-th distance is the core
  distance $\kappa_k$, a density estimate read off the data; \panel{d} keep
  the edges and you have a graph, which is all the rest of the toolbox needs.}
  \label{fig:knnsteps}
\end{figure}

\subsection{Cost, and the choice of $k$}

Exact $k$NN costs $O(n^2 d)$ for $n$ points in $d$ dimensions when computed
brute force, which for $3.6\times10^5$ field stars and $d = 16$ is about
$10^{12}$ distance evaluations --- feasible on a cluster, painful on a laptop.
The practical answers are spatial indices (KD-trees, ball trees, which degrade
badly above roughly 20 dimensions) and approximate graphs (NN-descent, the
algorithm UMAP uses, which is what makes it fast). This distinction matters
when reading performance claims: ``UMAP is fast'' is a statement about its
graph construction, not about exact $k$NN.

The choice of $k$ is the usual bias--variance trade: small $k$ gives a
high-variance density estimate that fragments clusters; large $k$ smooths away
small ones. Our pipeline uses $k = 15$ for both the UMAP graph and EVoC's
internal graph, and $k = 10$ for the purity score below --- all three are
conventional defaults, and all three are knobs a careful user should vary to
see whether the conclusion moves. \textbf{Section~\ref{sec:benchmark}} reports
one such variation: UMAP at $n_{\rm neighbours} = 5$ performs measurably worse
than at the default 15.

\subsection{Our use of it: kNN purity}

The one place where $k$NN enters our results directly is validation. For each
star that is a known member of a cluster, count the fraction of its 10 nearest
neighbours in the embedding that belong to the same cluster; average over
members. The result --- \emph{kNN purity} --- is a cohesion score with no
hyperparameters beyond $k$, no clustering step, and no way to tune yourself a
good answer:

\begin{equation}
\text{purity} = \frac{1}{|\mathcal{M}|}\sum_{i \in \mathcal{M}}
\frac{\left| \left\{ j \in \mathrm{kNN}_{10}(i) : \mathrm{cluster}(j) =
\mathrm{cluster}(i) \right\} \right|}{10},
\end{equation}

where $\mathcal{M}$ is the set of known members. It is the automated version
of the manual test in \citet{Kos:17} --- draw a polygon around the visual group
in the embedding and see whether the members are inside it --- and it is
reported for every method in \S\ref{sec:benchmark}. A purity near 1 means the
members of a cluster genuinely occupy a neighbourhood of the space; a purity
near the field fraction means the embedding knows nothing about them.

\begin{issues}[EXERCISES]
\begin{enumerate}
\item Derive the density estimate of \S 5.2 from the volume of a ball
  containing $k$ points. Then compute, for the DR19 abundance matrix, the
  spread of $\kappa_{15}$ across members and field stars. Is the cluster
  signal visible in the core distance alone, before any clustering?
\item Build the 15-NN graph of the member matrix and compute the purity of the
  known members in the \emph{raw} C-space. Compare with the purity after
  t-SNE, UMAP and EVoC. Which method moves the number most, and does it move
  it by separating members or by compressing the field?
\item ANN libraries advertise large speedups over exact $k$NN. Measure the
  disagreement rate between an approximate graph (UMAP's NN-descent) and an
  exact one on a random 5\,000-star subsample of the DR19 matrix. Is the
  approximation good enough that the clustering result is unchanged?
\end{enumerate}
\end{issues}

% 06 · DBSCAN — density, and the trouble with one threshold
\section{DBSCAN: FOLLOWING THE DENSITY}
\label{sec:dbscan}

\deckslides{slides 13--17, with slide 21 for the single-threshold problem.}

K-means asks how far a star is from a centre. DBSCAN --- density-based
spatial clustering of applications with noise \citep{Ester:96} --- asks a
different question: how many neighbours does a star have within a fixed
radius? That single change buys three things K-means cannot offer: clusters of
arbitrary shape, an explicit notion of noise, and no need to know $K$ in
advance.

\subsection{Definitions}

Fix two parameters: a radius $\varepsilon$ and a minimum count
\texttt{minPts}. Let $N_\varepsilon(\mathbf{x})$ be the set of points within
distance $\varepsilon$ of $\mathbf{x}$ (its $\varepsilon$-neighbourhood).

\begin{itemize}
\item \textbf{Core point.} $\mathbf{x}$ is a core point if
  $|N_\varepsilon(\mathbf{x})| \geq \texttt{minPts}$. Only core points can
  extend a cluster.
\item \textbf{Border point.} A non-core point that lies inside the
  $\varepsilon$-neighbourhood of a core point. It joins that cluster, but
  cannot recruit further points.
\item \textbf{Noise.} Everything else: not core, and not within $\varepsilon$
  of a core point. It is labelled $-1$ and never counted as a member of
  anything.
\item \textbf{Directly density-reachable.} $\mathbf{y}$ is directly
  density-reachable from $\mathbf{x}$ if $\mathbf{y} \in
  N_\varepsilon(\mathbf{x})$ and $\mathbf{x}$ is a core point. Note the
  asymmetry: a border point is reachable \emph{from} its core neighbour but
  not the other way round.
\item \textbf{Density-connected.} Two points are density-connected if there is
  a chain of core points, each directly reachable from the previous, touching
  both. Density-connectivity \emph{is} symmetric, and a DBSCAN cluster is a
  maximal set of density-connected points.
\end{itemize}

\begin{marginnote}[]
\entry{$\varepsilon$ (epsilon)}{The neighbourhood radius. It sets the density
threshold: a point is dense if it has \texttt{minPts} neighbours within
$\varepsilon$ of itself.}
\entry{minPts}{Minimum number of points (including the point itself) required
inside the $\varepsilon$-ball for a point to be a core point. Typical defaults
are $d+1$ for $d$ dimensions, or $\ln n$.}
\end{marginnote}

\begin{figure}[htbp]\centering
  \fig[0.94]{dbscan_film.png}
  \caption{DBSCAN growing a cluster, sampled at four points in the animation:
  core points are identified, a seed point starts a cluster, connected core
  points chain it along the crescent, and border points attach to the chain.
  The points that never enter a core's ball stay noise, permanently.}
  \label{fig:dbscanfim}
\end{figure}

\begin{figure}[htbp]\centering
  \begin{tabular}{@{}c@{\hskip6pt}c@{}}
    \fig[0.44]{dbscan_step1.png} &
    \fig[0.44]{dbscan_step2.png} \\[3pt]
    \fig[0.44]{dbscan_step3.png} &
    \fig[0.44]{dbscan_step4.png}
  \end{tabular}
  \caption{The algorithm in four panels: \panel{a} the two knobs, drawn as
  circles of radius $\varepsilon$ with a count of neighbours --- one point
  passes the \texttt{minPts} test, the other does not; \panel{b} every point
  classified as core (circles), border (triangles) or noise (crosses);
  \panel{c} core points within $\varepsilon$ of each other are chained
  (double-headed arrows), border points hang off the chain (single-headed);
  \panel{d} the result on two interleaved half-moons with uniform noise ---
  both crescents recovered whole, the scattered points discarded.}
  \label{fig:dbscansteps}
\end{figure}

\subsection{The algorithm}

The plain version is short enough to write out. For each point not yet
visited: if it is not a core point, mark it noise for now and move on; if it
is, start a new cluster and grow it by breadth-first search --- every core
point inside the current frontier's $\varepsilon$-neighbourhood joins the
cluster and becomes part of the frontier, and border points are attached as
they are met. The growth stops when the frontier is exhausted. The result does
not depend on the order in which points are visited, except for the assignment
of border points that are within $\varepsilon$ of two different clusters --- an
ambiguity the original paper leaves to the implementation, and a real
reproducibility hazard when comparing libraries.

Cost: with a spatial index, $O(n \log n)$ on average; without one, $O(n^2)$.
The clustering itself is single-pass and cheap; the work is in the
neighbourhood queries, which is the $k$NN machinery of \S\ref{sec:knn} again.

\subsection{Choosing the two knobs}

\texttt{minPts} is the easier of the two. It sets the order of the smallest
structure you are willing to call a cluster, and a common rule of thumb is
\texttt{minPts} $\approx d+1$ for $d$-dimensional data (here 17) or
$\ln n$ for large samples. Our pipeline uses \texttt{minPts} $=5$, which is
deliberately permissive: it keeps small clusters such as the Pleiades
candidates in play, at the cost of admitting noise-driven groups.

$\varepsilon$ is the harder one, and the standard data-driven recipe is a
$k$-distance plot: sort every point's distance to its $k$-th nearest
neighbour, plot the sorted values, and look for the knee --- the point where
the curve leaves the bulk of the data and starts to climb. The knee is where
the density threshold should sit. It is a heuristic, and it is exactly the
step that fails in this project:

\begin{issues}[THE SINGLE-THRESHOLD PROBLEM]
DBSCAN has one density threshold for the whole dataset. Our data have two
populations at very different densities: 357\,056 field stars forming a broad,
dense, chemically smooth background, and 25 clusters of 6--230 members each
sitting inside it. Set $\varepsilon$ for the clusters and the field becomes one
enormous cluster containing everything; set it for the field and the clusters
fragment into noise. There is no middle value that works, and the failure is
not a tuning problem --- it is the assumption of \textbf{Table~\ref{tab:assumptions}}
being false. The next two chapters exist because of it.
\end{issues}

The scale of the effect is visible in the results: on our DR19 run, an
un-normalised abundance space lets HDBSCAN merge the entire field into a
single group (recall $\approx 1.0$ with precision $\approx 0.001$), which is
the signature of a threshold set too high for the background. Row
normalisation (\S\ref{sec:data}) breaks that blob and is the single most
effective precision lever in the pipeline.

\begin{issues}[EXERCISES]
\begin{enumerate}
\item Draw the $\varepsilon$-ball classification by hand for the 2-D dataset
  in \textbf{Figure~\ref{fig:dbscansteps}} with $\varepsilon = 0.35$ and
  \texttt{minPts} $=4$, and mark every point core, border or noise. Then
  change \texttt{minPts} to 6 and say precisely which points change status.
\item Produce the $k$-distance plot for the DR19 member-plus-field matrix at
  $k=5$ and locate the knee. Repeat for $k=2$ and $k=20$. How does the knee
  move, and what does that say about whether a single $\varepsilon$ exists for
  these data?
\item Border points within $\varepsilon$ of two core points of different
  clusters are assigned by visit order in the reference implementation. Find
  two such points in a small synthetic example, and construct two visit orders
  that produce different partitions. Why does this matter when comparing
  scikit-learn with a GPU implementation?
\end{enumerate}
\end{issues}

% 07 · HDBSCAN* — density at every scale
\section{HDBSCAN$^*$: A HIERARCHY OVER EVERY DENSITY}
\label{sec:hdbscan}

\deckslides{slides 22--26 --- mutual reachability, the tree, and the long-lived branches.}

If the problem with DBSCAN is one global density threshold, the obvious repair
is to use all of them at once and let the data choose. That is what
HDBSCAN$^*$ does \citep{Campello:13, McInnes:17}: it builds a single
hierarchy of density levels, scores each candidate cluster by how long it
survives as the threshold moves, and returns the candidates that persist.
It is the clusterer attached to t-SNE and UMAP in our pipeline, and --- in a
different guise --- it is inside EVoC.

\subsection{From points to a density-aware metric}

The construction starts by making the metric itself density-aware. For each
point, compute the core distance $\kappa_k(\mathbf{x})$ of \S\ref{sec:knn}:
the distance to its $k$-th nearest neighbour. Then replace the original
distance between two points with the \emph{mutual reachability distance},

\begin{equation}
d_{\rm mreach}(\mathbf{a},\mathbf{b}) = \max\left\{\kappa_k(\mathbf{a}),\
\kappa_k(\mathbf{b}),\ d(\mathbf{a},\mathbf{b})\right\}.
\label{eq:mreach}
\end{equation}

Read it as a floor: two points in a dense region are pushed apart to at least
the local density scale, so a sparse point can never be ``close'' to anything.
Under this metric, a dense region becomes a compact object even when its
members are far apart in the original geometry --- which is exactly the
property DBSCAN lacked.

\begin{marginnote}[]
\entry{Mutual reachability}{The distance used by HDBSCAN$^*$: the maximum of
the two core distances and the true distance \textbf{(Equation~\ref{eq:mreach})}.
It makes dense regions self-similar across scales.}
\entry{Condensed tree}{The hierarchy after small clusters have been
progressively collapsed into their parents as the density threshold rises. Its
branches are the candidate clusters, weighted by lifetime (persistence).}
\end{marginnote}

\begin{figure}[htbp]\centering
  \begin{tabular}{@{}c@{\hskip8pt}c@{}}
    \fig[0.40]{mutual_reachability.png} &
    \fig[0.50]{knn_core_distance.png}
  \end{tabular}
  \caption{The ingredients of \textbf{Equation~\ref{eq:mreach}}, drawn in the
  deck. Left: two points, each ringed by its own core distance, with the
  straight-line distance between them; the mutual reachability distance is the
  largest of the three quantities. Right: the core distance itself --- the
  radius that contains exactly $k$ neighbours --- which is the local density
  estimate of \S\ref{sec:knn} written as a length.}
  \label{fig:mreach}
\end{figure}

\subsection{Building the hierarchy}

With the metric fixed, the rest is single-linkage clustering at every
threshold simultaneously, computed efficiently:

\begin{enumerate}
\item Build the minimum spanning tree (MST) of the mutual-reachability graph.
  One tree contains every single-linkage merge at every density level: the
  largest edge weight on the path between two points is the density at which
  they merge.
\item Sort the tree's edges by weight and cut them in that order, producing a
  dendrogram of nested partitions --- the density hierarchy
  (\textbf{Figure~\ref{fig:hdbscantree}}).
\item Condense: a cluster that splits into pieces, one of which is smaller
  than \texttt{min\_cluster\_size} and short-lived, is treated as noise rather
  than as a new cluster. This is what makes the tree robust against
  single-linkage's notorious chaining.
\item Score every remaining branch by \emph{stability}: the sum, over its
  members, of the difference between the density at which the member leaves
  the cluster and the density at which the cluster was born
  (\textbf{Equation~\ref{eq:stability}}). Selecting the set of branches with
  the largest total stability --- choosing a parent over its children when the
  parent scores higher --- gives the final labels.
\end{enumerate}

\begin{equation}
S(C_j) = \sum_{\mathbf{x} \in C_j} \left( \lambda_{\mathbf{x}} -
\lambda_{\rm birth}(C_j) \right),
\qquad \lambda = 1/d_{\rm mreach} .
\label{eq:stability}
\end{equation}

\begin{figure}[htbp]\centering
  \fig[0.72]{hdbscan_tree.png}
  \caption{The density hierarchy as a single-linkage dendrogram with a
  horizontal density cut. Everything above the line is one cluster; everything
  below it has split. HDBSCAN$^*$ does not commit to the line --- it scores
  every branch by how far it survives above its own birth level and keeps the
  branches with the best scores.}
  \label{fig:hdbscantree}
\end{figure}

\begin{figure}[htbp]\centering
  \fig[0.94]{hdbscan_film.png}
  \caption{The same data at four density levels, from the deck's animation:
  clusters appear as the threshold drops, some grow, and some merge into each
  other. HDBSCAN$^*$ sees the whole sequence at once and keeps the objects that
  live long enough to be structure rather than coincidence.}
  \label{fig:hdbscanfim}
\end{figure}

\begin{figure}[htbp]\centering
  \fig[0.62]{hdbscan_sweep.png}
  \caption{One density level of the sweep. Density-based clustering is not a
  fixed partition of the data but a family of partitions indexed by scale; the
  algorithm's job is to pick from the family without being told the scale in
  advance.}
  \label{fig:hdbscansweep}
\end{figure}

\subsection{What the knobs mean}

HDBSCAN$^*$ replaces DBSCAN's two knobs with two different ones, and both are
easier to reason about:

\begin{itemize}
\item \texttt{min\_cluster\_size} is a persistence floor, not a density
  threshold. It says: a structure smaller than this many points is not
  evidence of a cluster, so collapse it into its parent rather than report it.
  Our default is 5 --- deliberately low, because several of our real clusters
  are small --- and \S\ref{sec:benchmark} reports the effect of raising it to
  10, which on the DR17 data was the single largest improvement in recall.
\item \texttt{min\_samples} sets the $k$ inside the core distance of
  \textbf{Equation~\ref{eq:mreach}}: how many neighbours define the local
  density scale. Larger values give a smoother, more conservative density
  estimate --- the usual bias--variance trade of \S\ref{sec:knn}.
\item The metric is a third, silent knob: the whole construction presumes that
  a distance is meaningful. Cosine and Euclidean give different hierarchies;
  that is why \S\ref{sec:data} normalises rows, and why the fairness argument
  for comparing methods rests on the metric being shared.
\end{itemize}

\subsection{Why it fits our problem, and where it still fails}

The hierarchy is a natural fit for a field of 25 clusters of very different
sizes: it can, in principle, keep a 6-member cluster at one density and a
230-member cluster at another. In practice two things still bite.

First, \texttt{min\_cluster\_size} is a \emph{single} size floor. It cannot be
right for the Pleiades with 23 candidates and for M~67 with 230 members at the
same time --- one of them will be over-merged or the other fragmented. The
next chapter removes this knob entirely.

Second, and more fundamentally, the algorithm inherits the geometry it is
given. On our real abundances the residual failure is not the hierarchy but
the space: the field is a single dense blob that contains the clusters, and no
choice of \texttt{min\_cluster\_size} separates an object from its own
background. That is what motivates the embeddings of
\S\ref{sec:tsne}--\S\ref{sec:evoc}: change the geometry, then let HDBSCAN$^*$
work as designed.

\begin{issues}[EXERCISES]
\begin{enumerate}
\item For the four points $\mathbf{a},\mathbf{b},\mathbf{c},\mathbf{d}$ of a
  square grid, compute $d_{\rm mreach}$ with $k=2$ and verify that it is a
  valid metric. Which of the metric axioms does it break if you omit the
  $\kappa$ terms and use $\max\{d(\mathbf{a},\mathbf{b}),\ \text{const}\}$?
\item Compute the stability score of \textbf{Equation~\ref{eq:stability}} for
  a two-branch tree in which one branch contains 100 points that survive two
  decades of $\lambda$ and the other contains 10 points that survive five.
  Which does the algorithm prefer, and does that match your intuition about
  which is a better cluster?
\item Run HDBSCAN$^*$ on the DR19 member-plus-field matrix with row
  normalisation on and off. Record the number of clusters and the largest
  cluster fraction in each case, and identify which setting produces the
  ``one blob'' failure of \S\ref{sec:dbscan}.
\end{enumerate}
\end{issues}

% 08 · PLSCAN — persistence instead of a size floor
\section{PLSCAN: DROPPING THE SIZE FLOOR}
\label{sec:plscan}

\deckslides{slides 27--31 --- peaks, persistence, and the bars that survive.}

HDBSCAN$^*$ removed DBSCAN's density threshold but kept a single global
\texttt{min\_cluster\_size}: a point count, below which a structure is treated
as noise. That floor is a bad fit for our sample, where clusters range from 6
to 230 members. PLSCAN --- persistent multiscale density-based clustering,
\citet{Bot:25} --- takes the last step and removes it, replacing ``is this
group big enough?'' with ``how long does this group persist?''.

\subsection{Scale-space, and what persistence measures}

The idea comes from topological data analysis. Rather than committing to one
density threshold, consider the whole family of density estimates obtained by
smoothing the data at every scale. A real cluster appears as a \emph{peak}
that stays separate from its neighbours as the scale changes; a fluctuation
appears as a peak that is absorbed almost immediately. The range of scales
over which a peak exists as an independent object is its \emph{persistence},
and it is a statement about the data's structure rather than about a count.

The algorithm builds this structure explicitly:

\begin{enumerate}
\item \textbf{Scale-space.} Estimate the density field over a range of scales
  (equivalently: over a range of effective minimum cluster sizes), tracking the
  peaks at each scale.
\item \textbf{Leaf tree.} Record the peaks and their merges as a tree, where a
  \emph{leaf} is a density peak that exists in its own right --- the analogue
  of the branches of the HDBSCAN$^*$ condensed tree, but indexed by scale
  rather than by $k$-th-neighbour distance.
\item \textbf{Persistence trace.} For each leaf, measure the interval of scale
  over which the cluster survives as a separate object. Clusters by
  construction cannot be shorter-lived than their own persistence, which is
  what makes the criterion self-consistent.
\item \textbf{Selection.} Rate every cut of the leaf tree by the total
  persistence of the leaves it keeps, and take the cut with the highest total
  persistence; read off the labels at that cut.
\end{enumerate}

\begin{marginnote}[]
\entry{Persistence}{The range of minimum cluster sizes over which a cluster
exists as a separate structure, $s_{\max} - s_{\min}$ \citep{Bot:25}. A
long-lived peak is structure; a short-lived one is noise. It is measured in
stars --- the same units as the size floor it replaces --- which is why a
6-member cluster can outlive a 60-member fluctuation.}
\entry{Barcode}{One horizontal bar per cluster, its length equal to its
persistence, drawn on a common scale axis. Reading the barcode is the whole
selection step.}
\end{marginnote}

\begin{figure}[htbp]\centering
  \fig[0.94]{plscan_film.png}
  \caption{The persistence sweep, four frames of the deck's animation: the
  minimum-cluster-size axis is swept, and each peak's bar grows for as long as
  that peak remains a distinct cluster. Bars that stop early are fluctuations;
  bars that run the length of the sweep are structure.}
  \label{fig:plscanfim}
\end{figure}

\begin{figure}[htbp]\centering
  \begin{tabular}{@{}c@{\hskip5pt}c@{\hskip5pt}c@{}}
    \fig[0.30]{plscan_step1.png} &
    \fig[0.30]{plscan_step2.png} &
    \fig[0.30]{plscan_step3.png}
  \end{tabular}
  \caption{The three moves: \panel{a} find the density peaks; \panel{b} measure
  each peak's persistence and draw the barcode; \panel{c} keep the bars of the
  cut with the largest total persistence. The barcode replaces the size floor
  with a quantity measured over a range of sizes; the neighbour count $k$ that
  defines the core distance is still a choice, and when the trace has several
  peaks, so is which peak to report.}
  \label{fig:plscansteps}
\end{figure}

\begin{figure}[htbp]\centering
  \fig[0.78]{plscan_barcode.png}
  \caption{A three-peak density profile and its persistence barcode. The
  vertical extent of each bar is the range of ``minimum cluster size'' over
  which the corresponding cluster exists; the horizontal position identifies
  the peak. A min-cluster-size floor would cut all three bars at the same
  height, which is exactly the assumption PLSCAN avoids.}
  \label{fig:plscanbarcode}
\end{figure}

\subsection{What changes, and what does not}

Two honest statements, because the summary of this method is often oversold.

\textbf{What changes.} The criterion for ``is this a cluster?'' is no longer
a single size. A 6-member cluster that is a strong, long-lived density peak
survives; a 60-member fluctuation that dissolves as soon as the scale changes
does not, because the cut through the leaf tree is chosen by the trace over all
cuts rather than fixed in advance. That is the behaviour our sample needs, where
M~67 (230 members) and NGC~2158 (6 members) must be judged by the same rule.
What changes is not the units but who chooses: persistence is measured in stars,
and the data pick the value.

\textbf{What does not change.} The cut is still a choice --- PLSCAN makes it
automatically, by maximum total persistence, not without assumptions. The
neighbour count $k$ that sets the core distance is still an input, and their
own sensitivity study is the evidence that this matters less than HDBSCAN$^*$'s
size threshold. The assumption is moved to a quantity whose meaning is more
stable across datasets, and it is made explicit in the barcode, where a user can
see what a different cut would cost.
Nor does any of this repair a bad space: if the field and the clusters are not
separated in the geometry, no persistence criterion will separate them.

\subsection{Where it appears in this workbook}

PLSCAN enters our pipeline through EVoC, and that is worth understanding
before \S\ref{sec:evoc}. EVoC's final step sweeps the clustering of its
internal node embedding across scales, scores each resulting layer by the
total persistence of the clusters it contains, and returns the best-scoring
layer's labels. That is the PLSCAN idea used as a \emph{model-selection rule}
inside a larger algorithm --- the reason EVoC needs no
\texttt{min\_cluster\_size} from the user even though the density clustering
inside it would otherwise want one. On the deck's own terminology: the first
two columns of the EVoC pipeline come from UMAP's graph machinery, the middle
from HDBSCAN$^*$, and the last from PLSCAN.

\begin{issues}[EXERCISES]
\begin{enumerate}
\item Sketch a one-dimensional density with three peaks of different heights
  and widths. For each peak, mark the interval of smoothing scale over which
  it survives as a separate peak, and draw the resulting barcode. Now add a
  fourth peak that is real but tiny --- 2\% of the mass of the smallest other
  peak. Does the barcode distinguish it from a fluctuation?
\item In HDBSCAN$^*$ terms, what is the difference between ``this branch is
  shorter than \texttt{min\_cluster\_size}'' and ``this branch is shorter-lived
  than the others''? Construct a dataset where the two rules give different
  answers, and say which you would trust.
\item The persistence of a cluster is a range of minimum cluster sizes, i.e.\
  a number of stars (their Eq.~7). Using the member counts of
  \textbf{Table~\ref{tab:clusters}}, work out what a persistence of a given
  value implies for M~67 and for NGC~2158, and argue whether the same cut can
  serve both.
\end{enumerate}
\end{issues}

% 09 · Validation — how to know whether any of this is real
\section{VALIDATION: KNOWING WHEN A SCORE LIES}
\label{sec:validation}

\deckslides{slides 19--20 (internal and external validation) and slide 3.}

In a supervised problem, validation is arithmetic: hold out labels, count
mistakes. In clustering there is nothing to hold out, and the score you report
is chosen by you. This chapter is therefore the most important one in the
workbook, and the one in which we report the most unflattering numbers ---
including two of our own that forced us to withdraw a published claim.

\subsection{Internal indices: is the partition tidy?}

Three classical measures, all computed from the data and the partition alone.

\begin{itemize}
\item \textbf{Sum of squared errors.} The K-means objective of
  \textbf{Equation~\ref{eq:sse}} (\S\ref{sec:kmeans}).
  Monotone in the number of groups, so it is only useful as a curve against
  $K$, never as a score.
\item \textbf{Silhouette} \citep{Rousseeuw:87}. For a point $i$, let $a(i)$ be
  its mean distance to the other members of its own group and $b(i)$ its mean
  distance to the members of the nearest other group. Then
  $s(i) = (b(i)-a(i))/\max\{a(i),b(i)\} \in [-1,1]$, and the silhouette is the
  mean over points. It rewards compact, well-separated groups --- which is a
  statement about geometry, not about truth.
\item \textbf{Dip test} \citep{Hartigan:85}. A significance test for
  unimodality of a one-dimensional distribution. Its role here is upstream of
  clustering: it answers whether a variable (or a projection) shows any
  multipeaked structure at all. If nothing is multimodal, the clusters an
  algorithm returns are imposed by the algorithm.
\end{itemize}

All three share one weakness, and it is fatal if forgotten: \emph{a tidy
partition of unstructured data is easy to produce}. The silhouette of a
$K$-means solution on uniform noise is not zero, and the dip test is one
significance test among many you could run. Internal indices are for
diagnosing the behaviour of an algorithm, not for certifying a discovery.

\begin{marginnote}[]
\entry{Homogeneity}{How pure each predicted group is: 1 when every cluster an
algorithm reports contains a single true cluster, even if that true cluster is
split across several predictions.}
\entry{Completeness}{How intact each true cluster is: 1 when every true
cluster ends up in one predicted group, even if that group contains several
true clusters. Homogeneity and completeness trade off.}
\end{marginnote}

\subsection{External indices: comparison with something the algorithm never saw}

Here the labels come from outside the algorithm. Two families are used in this
workbook, and we use them for different tasks.

\textbf{Cluster-only separation} follows \citet{GarciaDias:19}. Take only the
known members, cluster them, and compare the predicted grouping with the true
one using an information-theoretic pair plus an assignment score:

\begin{align}
h &= 1 - \frac{H(C \mid K)}{H(C)}, &
c &= 1 - \frac{H(K \mid C)}{H(K)}, &
v &= \frac{2hc}{h+c},
\label{eq:merits}
\end{align}

where $C$ are the true cluster labels, $K$ the predicted ones, and $H$ the
Shannon entropy; accuracy is the fraction of correctly assigned stars after
matching predicted to true labels with the Hungarian algorithm (maximising the
diagonal of the confusion matrix). The pairing is not cosmetic. A method that
merges every cluster into one group has $c = 1$ and $h \approx 0$; a method
that shatters them has $h = 1$ and low $c$. Quoting one without the other is
how a broken result reads as a good one, and we do it ourselves in
\textbf{Table~\ref{tab:honest}} for exactly that reason.

\textbf{Field retrieval} follows the operational question: for each cluster,
take the predicted group that overlaps its true members most, and measure

\begin{equation}
\text{recall} = \frac{\text{true members recovered}}{\text{all true members}},
\qquad
\text{precision} = \frac{\text{true members in the group}}{\text{stars in the group}} .
\label{eq:rp}
\end{equation}

\textbf{Recovery fraction} \citep{Casamiquela:21} answers the question a
per-cluster study actually asks: how many of the known clusters came back? For
each true cluster, take the predicted group that holds most of its stars (the
best-overlap rule above). The cluster counts as \emph{recovered} at a threshold
$\theta$ when that group holds at least $\theta$ of the cluster \emph{and} at
least $\theta$ of the group is the cluster --- their per-cluster completeness and
homogeneity, our recall and precision, both at once:

\begin{equation}
\mathrm{RF}_{\theta} = \frac{\#\{\text{clusters recovered at } \theta\}}{\#\{\text{clusters}\}},
\qquad \theta = 0.40 \text{ or } 0.70 .
\label{eq:rf}
\end{equation}

Its companion is the \emph{statistical fraction}: the share of predicted groups
that recover no cluster at all, which is how the other paper measures the amount
of decoration in its own solution. We adopt this pair for the same reason we
adopt their task: of all the scores in this chapter it is the one whose chance
level is near zero whatever the sizes and number of clusters, so it survives
comparison between different samples. The implementation is
\texttt{cluster.baseline.recovery\_fraction}, and it reproduces their Table~2
exactly, which is the only external check on it available.

Two complements to \textbf{Equation~\ref{eq:rp}} appear in the results.
\emph{kNN purity} (\S\ref{sec:knn}) is parameter-free cohesion, immune to the
best-overlap protocol below. \emph{AUROC} answers a refusal question rather
than a discovery one: given a true member and a contaminant, how often does
the method rank the member as more plausible? It is used in
\S\ref{sec:benchmark} to show that the rejection power of spectra grows with
cluster age.

\begin{issues}[THE ``BEST OVERLAP'' IS A BEST CASE]
The protocol above scores a method by its single best-matching predicted
group per cluster. That is deliberately generous --- it is the only fair way to
compare algorithms that return different numbers of groups --- but it means a
method can look good by returning many groups, only one of which overlaps each
cluster. It also cannot be gamed into positive precision: the denominator is
the size of the predicted group, so a group with 3 true members in it and 300
field stars scores 0.01 however well it ``recovered'' those three.
\end{issues}

\subsection{Five ways a score lies}

Each of these is a case we hit. They are ordered from the oldest to the most
recent, and the fourth is the reason the validation chapter is in the middle
of the workbook rather than at the end.

\paragraph{1 · The degenerate partition.} A clustering that returns two groups
with two thirds of the stars in one of them produces a homogeneity of
$0.269$ --- a floor set by the collapse, not a measurement of the method. Two
of our PCA rows sat on that floor, and a first draft of this work described
the spectral latent as ``three times better than PCA''. It was three times
better than a crash. The check that catches it is a column nobody prints by
default: the fraction of stars in the largest predicted group. We now print
it, and \texttt{cluster head-to-head} flags degenerate rows automatically.

\paragraph{2 · Comparing different populations.} The abundance arm of our
first head-to-head was scored on 25 clusters and 1\,215 stars, while every
other arm was scored on 5 clusters and 55 stars. The abundance row therefore
looked worse than the others for a reason that had nothing to do with
abundances: it was being asked a harder question. The fix is structural, not
statistical --- intersect the sets of stars before scoring, so every row of
the table describes the same 982 stars and the same 25 clusters. Two tables in
this workbook have been rebuilt that way, and the rebuilt numbers are the ones
quoted.

\paragraph{3 · Row order, and duplicate stars.} Two checks on the same member
matrix, both of which should have been run before any of it was published.

\emph{Row order.} Permuting the rows of the abundance matrix --- a
transformation that leaves every pairwise distance exactly unchanged --- moves
the t-SNE homogeneity between $0.218$ and $0.560$, and the number of predicted
groups between 2 and 22. The sorted order gives $0.560$; the mean over eight
random orders is $0.268 \pm 0.100$. The geometry is identical in all nine
runs; what changes is the optimiser's path, which depends on initialisation
and on the order in which the Barnes--Hut approximation accumulates
interactions. Every other arm is stable to a few hundredths: UMAP on the same
abundances $0.579 \pm 0.009$, and the masked-AE latent $0.740 \pm 0.003$ under
t-SNE. The lesson is not ``t-SNE is broken'' but ``report the row-order spread
of any embedding with a stochastic initialisation'', because on small samples
it can exceed the effect you are claiming.

\emph{Duplicates.} An audit of the same matrix found that it contains repeat
entries for the same star: of 1\,002 member rows, 807 carry distinct
\texttt{APOGEE\_ID}s and 34 carry none, so 332 rows share an identifier with
another row (\texttt{scripts/population\_counts.py}). The copies are not identical --- the median spread between two
rows of the same star is $0.04$ dex per element, comparable to the abundance
differences we are trying to measure --- and they are concentrated: M~3
contributes 102 duplicate rows, M~67 ninety. This inflates the member counts
of \textbf{Table~\ref{tab:clusters}} (they are row counts, not star counts),
double-weights those clusters in every macro average, and is the most likely
explanation for the fragility of the abundances-only t-SNE row above. The
clean fix is to collapse each star's rows before clustering.

\paragraph{4 · Tuning on the score.} The lever sweep of
\S\ref{sec:benchmark} chose element weights by maximising the very metric it
then reported, on a single cluster at a single seed. That is a legitimate way
to \emph{find} a lever and an illegitimate way to \emph{quantify} it, and the
current run exposed the difference: the same switch that lifted t-SNE recall
on the DR17 data ($0.09 \rightarrow 0.44$) lowers precision on the DR19 data
($0.83 \rightarrow 0.24$). A lever is a hypothesis; only a re-run on the
current data, quoted with both recall and precision, turns it into a result.

\paragraph{5 \textperiodcentered\ The chance level of a mutual-information score moves
with the sample.} Small groups inflate $h$, $c$ and $V$. Shuffle
\citet{Casamiquela:21}'s 175 stars at random into groups with the sizes they
observe --- 31 clusters, most of them tiny --- and the V-measure is $0.51$; do the
same on our 982 stars in 25 clusters and it is $0.06$--$0.12$. Their published
$V = 0.55$ therefore sits $0.04$ above its null, and our $V = 0.58$ sits about
$0.47$ above ours. The two numbers are not comparable, and a table that prints
them side by side without their nulls asserts a comparison that does not exist.
\textbf{Equation~\ref{eq:rf}} does not have this problem: chance level $0.04$
for their sample shape and $\le 0.001$ for ours. This is the reason
\S\ref{sec:literature} is written in recovery fractions.

\subsection{The protocol this workbook follows}

Out of all that, six rules, which the numbers in
\S\ref{sec:benchmark}--\S\ref{sec:physics} obey:

\begin{enumerate}
\item \textbf{Same population or no comparison.} Arms are intersected on the
  stars they share before any score is computed.
\item \textbf{Mean $\pm$ standard deviation over seeds.} Seven seeds
  ($42, 0, 1, 2, 7, 13, 99$), never a single run, and never a best-of.
\item \textbf{Kinematics are the ceiling, not a competitor.} They define the
  labels; the interesting question is always how close chemistry or spectra
  get to that ceiling.
\item \textbf{Parameter-free scores alongside tuned ones.} kNN purity and
  AUROC cannot be improved by tuning, which makes them the sanity check on
  everything else.
\item \textbf{Report the degenerate rows.} Largest-group fraction and number
  of predicted groups are printed next to every score, so a collapse is
  visible without recomputation.
\item \textbf{Recompute the chance level before comparing with a published
  score.} A score quoted without its null is a number without units, and the
  null moves when the sample does (rule 5 above).
\end{enumerate}

\begin{figure}[htbp]\centering
  \begin{tabular}{@{}c@{\hskip8pt}c@{}}
    \fig[0.44]{confusion_tsne_chem.png} &
    \fig[0.44]{confusion_tsne_kin.png}
  \end{tabular}
  \caption{The confusion matrix of a t-SNE clustering of the members, with the
  true cluster on one axis and the predicted group on the other
  (\textbf{Equation~\ref{eq:merits}} visualised). Left: abundances only --- one
  predicted group absorbs most of the true clusters, which is the
  completeness-rich, homogeneity-poor blob described above. Right: adding
  kinematics. The blocks along the diagonal are what ``the clusters are
  separated'' looks like in a matrix, and the difference between the two panels
  is the central result of \S\ref{sec:benchmark}.}
  \label{fig:confusion}
\end{figure}

\begin{figure}[htbp]\centering
  \fig[0.78]{proper_motions.png}
  \caption{Why kinematics can referee. Each panel shows proper motion in right
  ascension and declination for the stars around one cluster; the tight group
  sharing a velocity is the cluster, the surrounding cloud is the field. These
  are the observables that define the labels, so they can never be an input to
  a method that is scored against them.}
  \label{fig:propermotions}
\end{figure}

\begin{issues}[EXERCISES]
\begin{enumerate}
\item Show that a clustering that puts every star in a single group has
  $h = 0$ and $c = 1$, and that a clustering that gives every star its own
  group has $h = 1$ and $c \approx 0$. Where does accuracy sit in each limit?
\item Reproduce the row-order check: cluster the DR19 member matrix with
  t-SNE$\,+\,$HDBSCAN nine times, once sorted and eight times with a random
  row order, fixing the seed. Report the spread of homogeneity and of the
  number of predicted groups. Now do the same with the masked-AE latent. Why
  is one stable and the other not?
\item Collapse the duplicate rows of the member matrix (for example by taking
  the median abundance vector per \texttt{APOGEE\_ID}) and re-run the
  cluster-only table. Which rows of \textbf{Table~\ref{tab:honest}} move, and
  does M~3's share of the sample change the conclusion of the ablation?
\item Design a $\sigma$-clipped chemical membership and use it to score a
  kinematic clustering. Based on \S\ref{sec:data}, explain in three sentences
  why the resulting number would be meaningless --- and what the equivalent
  mistake would look like in your own field.
\end{enumerate}
\end{issues}

% 10 · t-SNE — neighbour probabilities and the crowding problem
\section{t-SNE: A MAP OF NEIGHBOURHOODS}
\label{sec:tsne}

\deckslides{slides 32--38 --- the two probability models and the descent between them.}

Everything so far has treated the 16-dimensional abundance vector as a space
in which distance means something. That assumption is weaker than it looks. In
high dimensions, distances concentrate: the ratio of the farthest to the
nearest neighbour of a point tends to 1, so ``close'' and ``far'' describe
almost the same thing and a density-based method loses contrast. t-SNE
\citep{vanderMaaten:08}, building on stochastic neighbour embedding
\citep{Hinton:02}, answers this by giving up on preserving distances at all
and preserving \emph{neighbourhoods} instead --- then drawing the result in two
dimensions.

\begin{figure}[htbp]\centering
  \begin{tabular}{@{}c@{\hskip5pt}c@{\hskip5pt}c@{}}
    \fig[0.30]{tsne_step1.png} &
    \fig[0.30]{tsne_step2.png} &
    \fig[0.30]{tsne_step3.png}
  \end{tabular}
  \caption{The three moves of t-SNE, from the deck: \panel{a} in the high-D
  space, convert distances to neighbour probabilities (here the point's
  neighbourhoods are drawn as spheres of decreasing probability); \panel{b} in
  the 2-D map, do the same conversion with a heavy-tailed kernel; \panel{c}
  move the points downhill on the Kullback--Leibler divergence until the two
  sets of probabilities agree as closely as they can.}
  \label{fig:tsnesteps}
\end{figure}

\subsection{From distances to probabilities}

For each point $i$, t-SNE defines the probability that $j$ is its neighbour in
the high-dimensional space as a Gaussian in the distance,

\begin{equation}
p_{j \mid i} = \frac{\exp\!\left(-\|\mathbf{x}_i -
\mathbf{x}_j\|^2 / 2\sigma_i^2\right)}
{\sum_{k \neq i} \exp\!\left(-\|\mathbf{x}_i - \mathbf{x}_k\|^2 /
2\sigma_i^2\right)} ,
\label{eq:phigh}
\end{equation}

and symmetrises, $p_{ij} = (p_{j|i} + p_{i|j})/2n$. The novelty is $\sigma_i$:
each point gets its \emph{own} bandwidth, chosen so that the neighbourhood
distribution has a fixed entropy. The user-facing parameter is that entropy,
expressed as \emph{perplexity},

\begin{equation}
\mathrm{Perp}(P_i) = 2^{H(P_i)}, \qquad
H(P_i) = -\sum_j p_{j|i} \log_2 p_{j|i},
\label{eq:perplexity}
\end{equation}

found by binary search on $\sigma_i$. Perplexity is the effective number of
neighbours each point is required to have, and it is the parameter you must
justify. \citet{vanderMaaten:08} suggest 5--50; a widely used default is
$\sqrt{n}$ or $n/100$; our pipeline's default is 30, and the swap to 50
measurably improved the layout for our $\sim\!5\,000$-star samples.

\begin{marginnote}[]
\entry{Perplexity}{The effective number of neighbours each point is given in
t-SNE's high-dimensional probability distribution
(\textbf{Equation~\ref{eq:perplexity}}). It is the only required
hyperparameter, and it is a statement about the scale of structure you are
looking for.}
\entry{Crowding problem}{In a low-dimensional map there is not enough room to
place moderately distant points at moderate distances. t-SNE's heavy-tailed
low-dimensional kernel fixes it by making the map's tail fat.}
\end{marginnote}

\subsection{The low-dimensional half, and the crowding problem}

The map is a set of 2-D positions $\mathbf{y}_i$. Their neighbour
probabilities use a Student-$t$ kernel with one degree of freedom,

\begin{equation}
q_{ij} = \frac{\left(1 + \|\mathbf{y}_i - \mathbf{y}_j\|^2\right)^{-1}}
{\sum_{k \neq l} \left(1 + \|\mathbf{y}_k - \mathbf{y}_l\|^2\right)^{-1}} ,
\label{eq:qlow}
\end{equation}

and the layout is found by minimising the Kullback--Leibler divergence
$D_{\mathrm{KL}}(P \,\|\, Q)$ by gradient descent. The heavy tail in
\textbf{Equation~\ref{eq:qlow}} is the whole point. With a Gaussian map
kernel, the volume available at radius $r$ in two dimensions grows only as
$r$, while the volume of a shell in 16 dimensions grows as $r^{15}$: a point
that should be moderately distant from its neighbours has nowhere to go, and
the optimiser either piles dissimilar points on top of each other or flings
them arbitrarily far away. This is the crowding problem, and the $t_1$ kernel
solves it by allowing large distances to be cheap: $q_{ij} \sim
1/\|\mathbf{y}_i - \mathbf{y}_j\|^2$, so a point can sit far from another at
almost no cost in divergence. The resulting gradient has an intuitive form ---
each point is pulled towards its high-D neighbours and pushed away from the
others, with the repulsion weak at long range --- and it is what
\textbf{Figure~\ref{fig:tsnefilm}} animates.

\begin{figure}[htbp]\centering
  \fig[0.94]{tsne_film.png}
  \caption{The descent, four frames of the deck's animation: from the initial
  (PCA) layout, through the early exaggeration phase in which clusters
  coalesce while empty space opens up between them, to the settled map. The
  exaggeration phase is a device, not physics: it deliberately over-states
  attraction at the start to escape a poor initial arrangement.}
  \label{fig:tsnefilm}
\end{figure}

\begin{issues}[HOW TO READ A t-SNE MAP]
The picture is a map of neighbour probabilities, and the following are all
\emph{not} readable from it. (i) Distances between clusters: only the relative
order of neighbourhoods is preserved, so a gap can appear between two groups
whose members are mutual neighbours. (ii) Cluster sizes and densities: t-SNE
normalises local density away, so a tight blob and a loose one can have the
same number of members. (iii) Global geometry: rotation, reflection and the
arrangement of clusters are arbitrary. (iv) The absence of a gap: a method
that must place every point somewhere will place well-separated points on the
same island if the perplexity is wrong.
\end{issues}

\subsection{What t-SNE does not do}

It does not cluster. The output is a 2-D embedding, and the visual groupings
are an artefact of a continuous optimisation. The workflow that has become
standard --- and that our pipeline uses --- is to treat the embedding as a
feature space and cluster \emph{that} with a density-based method, normally
HDBSCAN$^*$ (\S\ref{sec:hdbscan}) or EVoC (\S\ref{sec:evoc}). The
distinction matters when quoting results: a t-SNE row in our tables is
really ``t-SNE followed by HDBSCAN$^*$'', and both steps contribute to the
score.

Three further practical facts, each of which we have hit:

\begin{itemize}
\item \textbf{Initialisation changes the answer.} The original implementation
  started from a random layout; the modern default starts from a PCA of the
  input. Both are stochastic in practice, and on our abundance matrix the
  difference is not subtle: see the row-order audit of \S\ref{sec:validation},
  where the same data gave homogeneity between 0.218 and 0.560 depending on
  the order of the input rows.
\item \textbf{The distance metric is fixed to Euclidean.} t-SNE has no
  metric parameter, which is why \S\ref{sec:data}'s decision to make
  Euclidean distance and cosine similarity equivalent on our matrix is not a
  detail: it is the only reason the t-SNE, UMAP and EVoC arms are comparable
  at all.
\item \textbf{Cost grows faster than linearly.} The original $O(n^2)$
  formulation is unusable past a few thousand stars; the Barnes--Hut
  approximation \citep{vanderMaaten:14} brings it to
  $O(n \log n)$ per iteration, at the cost of the order dependence noted
  above. Our field runs embed 30\,107 stars in this regime.
\end{itemize}

\subsection{t-SNE on real clusters}

The published anchor for this section is \citet{Kos:17}, who applied t-SNE to
GALAH abundances --- 13 elements, calibrated with The Cannon --- and recovered 7
of the 9 clusters in their sample, the Pleiades among them, as coherent groups in
a two-dimensional map. \textbf{Figure~\ref{fig:kostse}} reproduces their figure,
and it is worth being precise about what it shows: a known cluster, its members
in red, sitting inside a map built from the surrounding field. It is a positive
result, and it is also a reminder of the asymmetry in this literature: the
clusters were selected because they were already known, and the method is judged
on whether it finds them again.

Our own cluster-only results for the abundance matrix illustrate the failure
mode the caveats above predict, and they do it in two registers. The single run
of \textbf{Table~\ref{tab:baseline}} gives the t-SNE arm the highest completeness
of the three methods ($0.879$) and the lowest homogeneity ($0.292$), the
signature of a map in which most clusters have been drawn into one region. The
seven-seed average of \textbf{Table~\ref{tab:honest}} looks better ($0.560$
homogeneity, $0.601$ completeness) --- and \S\ref{sec:validation} shows that the
same matrix moves between $0.218$ and $0.560$ on homogeneity, and between 2 and
22 groups, when only the order of the input rows changes. The neighbourhoods are
preserved, and the neighbourhoods were not separated to begin with.

\begin{figure}[htbp]\centering
  \fig[0.62]{kos_pleiades_tsne.png}
  \caption{t-SNE of GALAH abundances around the Pleiades, as published by
  \citet{Kos:17} and used in the lecture: the red points are stars known to be
  cluster members, and the inset zooms in on where they concentrate, marking the
  two groups the method finds (A, B) and the two member candidates it adds
  (1, 2). A known cluster, recovered from field data --- the benchmark that
  strong chemical tagging has to beat.}
  \label{fig:kostse}
\end{figure}

\begin{issues}[EXERCISES]
\begin{enumerate}
\item Implement the binary search for $\sigma_i$ in
  \textbf{Equation~\ref{eq:perplexity}} that fixes perplexity at a target
  value for every point. Plot the resulting $\sigma_i$ against local density
  on the DR19 matrix. Does the adaptive bandwidth do what it is supposed to
  do?
\item Embed the member matrix with perplexities 5, 30 and 100, and for each
  compute the kNN purity of \S\ref{sec:knn} plus the cluster-only homogeneity.
  Which perplexity would you choose, on which evidence, and does the choice
  change the paper's conclusion?
\item Take the t-SNE embedding of the abundances and run HDBSCAN$^*$ with
  \texttt{min\_cluster\_size} = 5, 15 and 50. Report the number of groups and
  the largest group fraction each time. Identify which run is degenerate and
  say what the degeneracy looks like in the map.
\item The low-dimensional kernel is $q_{ij} \propto (1 +
  \|\mathbf{y}_i-\mathbf{y}_j\|^2)^{-1}$. Replace it with a Student-$t$ kernel
  of $\nu = 5$ degrees of freedom, re-derive the gradient, and explain what
  you expect to change in the map's long-range behaviour.
\end{enumerate}
\end{issues}

% 11 · UMAP — fuzzy graphs and a cross-entropy layout
\section{UMAP: A GRAPH, THEN A LAYOUT}
\label{sec:umap}

\deckslides{slides 39--42 --- graph, fuzzy edges, layout.}

t-SNE converts a fully connected distance matrix into probabilities and then
optimises a two-dimensional layout against them. UMAP \citep{McInnes:18} takes
a different route to a similar-looking picture: build a weighted
nearest-neighbour graph, decide what that graph's low-dimensional analogue
should look like, and find positions that make the two agree. The construction
is more explicit about its assumptions, it is faster, and --- as
\S\ref{sec:benchmark} shows --- it behaves differently on our data in a way
that teaches something about both methods.

\begin{figure}[htbp]\centering
  \begin{tabular}{@{}c@{\hskip5pt}c@{\hskip5pt}c@{}}
    \fig[0.30]{umap_step1.png} &
    \fig[0.30]{umap_step2.png} &
    \fig[0.30]{umap_step3.png}
  \end{tabular}
  \caption{UMAP in three moves, from the deck: \panel{a} build the
  $k$-nearest-neighbour graph of the high-dimensional data; \panel{b} weight
  its edges by local distance and treat the result as a fuzzy set --- every
  edge has a probability of existing; \panel{c} lay the graph out in low
  dimensions so that strong edges are short and the layout stays spread out,
  by minimising the cross-entropy between the two edge distributions.}
  \label{fig:umapsteps}
\end{figure}

\subsection{The high-dimensional graph}

For each point $i$, let $\rho_i$ be the distance to its nearest neighbour ---
a local connectivity term that makes the metric robust in dense regions ---
and let $\sigma_i$ be a bandwidth chosen so that the total edge weight out of
$i$ is a fixed constant:

\begin{equation}
w_{ij} = \exp\!\left(-\frac{d(i,j) - \rho_i}{\sigma_i}\right).
\label{eq:umapw}
\end{equation}

Small distances between close neighbours are saturated at $\exp(0)=1$; the
bandwidth is solved per point by binary search, exactly as t-SNE solves for
$\sigma_i$ in \textbf{Equation~\ref{eq:phigh}}, with the target here being
$\sum_j w_{ij} = \log_2 k$ rather than a perplexity. The result is a directed
graph, made symmetric by the fuzzy-set union
$w_{ij} + w_{ji} - w_{ij}w_{ji}$.

Two design decisions are worth naming. First, the graph is built from an
approximate $k$NN search (NN-descent, \S\ref{sec:knn}), which is why UMAP
scales to $10^6$ points where t-SNE does not. Second, the graph is
\emph{local}: edges exist only to the $k$ nearest neighbours, so UMAP never
sees a long-range distance. Every claim about UMAP preserving ``more global
structure'' than t-SNE has to be read against that fact --- the long ranges it
appears to preserve are produced by the layout's initialisation and by the
repulsion term, not by information in the graph.

\subsection{The low-dimensional half}

In the map, the probability assigned to an edge is a parameterised curve

\begin{equation}
q_{ij} = \left(1 + a\,\|\mathbf{y}_i - \mathbf{y}_j\|^{2b}\right)^{-1},
\label{eq:umapq}
\end{equation}

with $a$ and $b$ fitted so that the curve approximates the shape of the
high-dimensional weights for the requested \texttt{min\_dist}. The objective is
the cross-entropy between the two edge distributions, minimised by stochastic
gradient descent with negative sampling --- a handful of random non-edges
attracted or repelled per real edge, which is the second reason UMAP is fast.

\begin{marginnote}[]
\entry{\texttt{n\_neighbors}}{The size of the local graph, and UMAP's analogue
of perplexity: it sets the scale of structure the layout can see. Small values
concentrate on very local structure; large values give a more global, blobier
picture.}
\entry{\texttt{min\_dist}}{How tightly points may pack in the layout. It
changes the appearance of density, not the structure found: a low value makes
clusters look like tight clumps.}
\end{marginnote}

\subsection{Parameters, and why our defaults are what they are}

UMAP exposes more knobs than t-SNE, and three of them matter here.

\begin{itemize}
\item \texttt{n\_neighbors} (default 15 in our pipeline, 15 in the library).
  On the DR17 lever sweep this was the difference between a working and a
  broken layout: at \texttt{n\_neighbors} $=5$ the recall on M~67 fell from
  $0.415$ to $0.071$, while at $30$ it was indistinguishable from the default
  ($0.410$). The lesson generalises: too few neighbours gives a graph that
  fragments the field, not one that isolates the cluster.
\item \texttt{metric}. Our matrix is L2-normalised (\S\ref{sec:data}), so
  Euclidean and cosine distances are monotonically related and the metric
  should be nearly irrelevant. In the same sweep, forcing
  \texttt{metric=cosine} cut recall from $0.415$ to $0.044$. The nominal
  geometry did not change; the implementation's internal handling of the graph
  did. Report the setting you used.
\item \texttt{random\_state}. UMAP's optimisation is genuinely stochastic ---
  unlike the t-SNE rows in our tables, which came out identical across the
  seven seeds we tried, UMAP's homogeneity spread is $\pm0.011$ on the
  masked-AE latent. That is small, but it is the same order as some of the
  gaps quoted between arms.
\end{itemize}

\begin{figure}[htbp]\centering
  \fig[0.94]{umap_film.png}
  \caption{The layout being optimised, four frames of the deck's animation:
  the graph is embedded with attraction along the edges and repulsion between
  sampled non-neighbours, and the structure emerges as the cross-entropy
  falls. There is no phase analogous to t-SNE's early exaggeration; the layout
  is driven the whole way by the same objective.}
  \label{fig:umapfilm}
\end{figure}

\subsection{How it differs from t-SNE in practice}

The two methods answer the same question --- produce a low-dimensional map that
keeps neighbours together --- and disagree in ways that matter for the rest of
this workbook:

\begin{itemize}
\item \textbf{Speed and scaling.} UMAP is the only arm we can run on the full
  357\,056-star field without sampling; the density-clustering arm of our
  pipeline therefore defaults to it for field work, and the t-SNE arm is run in
  region mode or on subsamples.
\item \textbf{What ``noise'' looks like.} UMAP has no noise model. Everything
  gets a position, so a cluster in a UMAP map is a density peak and nothing
  more --- which is why the pipeline clusters the embedding afterwards rather
  than reading groups off the picture.
\item \textbf{Blob formation.} On our abundance matrix, UMAP reaches a higher
  cluster-only homogeneity than t-SNE ($0.578$ vs $0.560$) while producing
  larger, fewer groups --- and in the field-retrieval task it recovers all
  978 members at a precision of $0.0013$, that is, by declaring a large part of
  the field to be cluster members. \textbf{Table~\ref{tab:field}} is the
  clearest illustration in the workbook of why recall and precision must be
  read as a pair, and of why the ``which method is better'' question has no
  answer until the task is specified.
\end{itemize}

\begin{issues}[EXERCISES]
\begin{enumerate}
\item Implement the per-point binary search for $\sigma_i$ in
  \textbf{Equation~\ref{eq:umapw}} that fixes $\sum_j w_{ij} = \log_2 k$, and
  verify numerically on a random dataset. Where does the local-connectivity
  term $\rho_i$ matter most: dense regions or sparse ones?
\item For a small dataset, compute UMAP's edge weights and then run the
  layout twice: once with PCA initialisation and once with a random one. Do
  the two layouts agree on which points are neighbours? Which of them would
  you publish, and what would you have to state about it?
\item \texttt{min\_dist} changes the appearance of the map without changing
  the graph. Using a fixed embedding, vary \texttt{min\_dist} over
  $\{0.0, 0.1, 0.5, 1.0\}$ and compute the cluster-only homogeneity each time.
  Does the clustering result move, and does that support or undercut the use
  of UMAP maps as evidence?
\item UMAP's repulsion term acts on sampled non-edges, so its notion of
  ``far'' is an artefact of negative sampling. Construct a small dataset in
  which the layout places two true clusters far apart and the reverse, by
  changing only \texttt{random\_state}. What does that imply for reading
  distances from a published UMAP figure?
\end{enumerate}
\end{issues}

% 12 · EVoC: embedding, density, and persistence in one object
\section{EVoC: ONE FIT, THREE STAGES}
\label{sec:evoc}

\deckslides{slides 43--45: the fusion, the pipeline, and why C-space.}

The three methods so far are composed by hand: t-SNE or UMAP produces an
embedding, then HDBSCAN$^*$ (or a persistence-driven variant of it) produces
labels, then the user decides how the two are coupled. EVoC
\citep{EVoC} removes the seam. It takes the graph, embeds it, clusters it, and
chooses the clustering itself, in a single fit, and it is the method our
pipeline uses for its production runs.

Read the deck's pipeline figure first: the point of the method is that the
three stages share one objective rather than being tuned against each other.

\begin{figure}[htbp]\centering
  \fig[0.92]{evoc_pipeline.png}
  \caption{The EVoC pipeline. From left to right: a nearest-neighbour graph is
  built from the raw data with a native metric, the graph is embedded into a
  low-dimensional space by a graph-embedding objective, and the embedded points
  are clustered with a density-based method whose candidate clusterings are
  scored by persistence. The output is labels, not a picture.}
  \label{fig:evocpipeline}
\end{figure}

\subsection{Stage 1: the graph, built on the metric you meant}

EVoC's first stage is a $k$-nearest-neighbour graph (\S\ref{sec:knn})
constructed directly from the input representation, with a configurable
metric. This matters in our setting for a specific reason: it is the one arm
that can use the \emph{raw} L2-normalised abundance matrix and its native
cosine geometry, without an intervening Euclidean-only embedding. Since
\S\ref{sec:data} normalises every row to unit length, Euclidean and cosine
distances are monotonically related on our matrix, and EVoC's use of cosine is
a deliberate choice rather than an accident: in a compositional space, only
the ratios between elements are meaningful, and a vector's length is an
artefact of the abundance sum.

\begin{figure}[htbp]\centering
  \begin{tabular}{@{}c@{\hskip6pt}c@{}}
    \fig[0.44]{evoc_step1.png} &
    \fig[0.44]{evoc_step2.png} \\[3pt]
    \fig[0.44]{evoc_step3.png} &
    \fig[0.44]{evoc_step4.png}
  \end{tabular}
  \caption{The four stages as the lecture builds them: \panel{a} the
  $k$-nearest-neighbour graph of the data; \panel{b} an embedding of the graph
  nodes (the graph's structure is preserved, not its geometry); \panel{c}
  density-based clustering of the embedded points; \panel{d} the persistence
  score used to pick which clustering layer to return.}
  \label{fig:evocsteps}
\end{figure}

\subsection{Stage 2: an embedding for clustering, not for looking at}

The embedding stage optimises a graph-embedding objective: adjacent nodes
close, non-adjacent nodes apart, with the dimension chosen by the algorithm
rather than by the user. It is worth being explicit about what this is
\emph{not}. It is not a visualisation; there is no meaningful scatter plot at
the end of it, and we never draw one. Its purpose is to place the graph in a
space where a density-based clusterer can separate its communities, which is a
different objective from t-SNE's or UMAP's two-dimensional layout problem.

This is also where the most common confusion in the literature arises. When a
paper says ``UMAP followed by HDBSCAN$^*$'', there are two algorithms and
therefore two independent sources of error; when a paper says ``EVoC'', the
cluster structure is chosen inside the fit, by persistence, from the same graph
that the embedding was built on. The distinction is not about quality but about
where the arbitrary decisions live, and it is why our benchmark reports all
three, since each one answers a different question about the same data.

\subsection{Stage 3: layers, persistence, and no \texttt{min\_cluster\_size}}

The clustering stage produces not one answer but a family of them, indexed by
scale. Each candidate layer is scored by the total persistence of the clusters
it contains (the PLSCAN criterion of \S\ref{sec:plscan}) and the best-scoring
layer's labels are returned. Consequences that matter in practice:

\begin{itemize}
\item \textbf{No size floor to set.} The user is never asked for
  \texttt{min\_cluster\_size}, which is the single knob that most changed
  results in our DR17 sweeps (\S\ref{sec:benchmark}: raising it from 5 to 10
  moved t-SNE recall from $0.087$ to $0.404$). EVoC replaces it with an
  internal persistence criterion, which is more stable across datasets but
  less inspectable: you cannot see which layers were considered.
\item \textbf{The number of clusters is an output.} On the field runs, EVoC
returns a few dozen groups on average, none of them large: one of the reasons
its recall is high and its precision catastrophic
(\textbf{Table~\ref{tab:field}}).
\item \textbf{The returned layer is a model-selection decision.} Persistence
  is a principled criterion, not a guarantee: if the field and the clusters
  are not separated in the graph's geometry, the best-scoring layer is the one
  that splits the field most persistently.
\end{itemize}

\subsection{What EVoC does for us, and what it does not}

Positive first. EVoC is the only arm that needs no post-hoc clustering choice,
it works directly on the cosine geometry that our data actually have, it is
the cheapest of the three to run on large samples, and on cluster-only
separation it is competitive: homogeneity $0.420 \pm 0.043$ on abundances,
$0.733 \pm 0.008$ on the 64-dimensional PCA latent, both measured over seven
seeds on the same 982 stars as the other arms.

Now the part that took us a while to accept. Two results in the literature
around this method, and two of ours, all point the same way:

\begin{issues}[WHAT WE HAVE TO SAY AGAINST OUR OWN FAVOURITE] \textbf{The seed
spread is larger than several of the gaps.} EVoC's homogeneity on abundances is
$0.420 \pm 0.043$ over seven seeds: a spread that overlaps with the t-SNE and
UMAP arms. In the five-cluster head-to-head the spread is $\pm0.08$--$0.11$.
Any ranking that depends on differences of that size is not a ranking.

\textbf{The self-supervision claim is a t-SNE/UMAP result, not an EVoC result.}
On the corrected same-population table, the masked autoencoder scores $0.70 \pm
0.08$ against $0.64 \pm 0.08$ for abundances on EVoC (overlapping within
errors) while the \emph{supervised} CNN wins EVoC outright at $0.80 \pm 0.03$
on the four-cluster sample. We withdrew the sentence that said the
abundance-free latent wins everywhere.

\textbf{High recall at near-zero precision.} On the all-sky field-retrieval
task, EVoC recovers $0.557$ of the true members at a precision of $0.0033$:
the groups it returns are mostly field stars, by two orders of magnitude. It
is the same failure that UMAP shows in that table, and it is the reason the
pipeline's field work runs a rejection stage after the clustering rather than
trusting the labels.
\end{issues}

\begin{issues}[EXERCISES]
\begin{enumerate}
\item EVoC clusters an embedding built from the graph, and HDBSCAN$^*$
  clusters the embedding built by t-SNE or UMAP. Explain, in three sentences
  each, why the first is more self-consistent and why the second is easier to
  diagnose when it fails.
\item On the member matrix, run EVoC with the seed varied over the same seven
  values used in this workbook, and report homogeneity as mean $\pm$ standard
  deviation. Now do the same for t-SNE with a fixed seed but nine row orders
  (\S\ref{sec:validation}). Which source of variability is larger, and what
  does that say about which number belongs in a paper?
\item Take the EVoC labels on the all-sky field run and compute, for each
  returned group, the fraction of its members that are known field stars.
  Plot that distribution. How many groups would you have to ignore before the
  remaining ones have precision above 0.5?
\item The persistence criterion selects a clustering layer without user
  input. Construct a two-dimensional dataset with a dense uniform background
  and one small dense cluster, and report which layer EVoC selects. Does the
  criterion find the cluster, the background, or something else?
\end{enumerate}
\end{issues}

% 13 · The benchmark — what we measured, and what it does not support
\section{THE BENCHMARK: WHAT WE MEASURED}
\label{sec:benchmark}

\deckslides{slides 46--47 for the benchmark and slide 48 for the lessons in short form.}

This chapter is the case study: the methods of
\S\ref{sec:kmeans}--\S\ref{sec:evoc} run on the data of \S\ref{sec:data},
scored by the protocol of \S\ref{sec:validation}. The order of business is the
benchmark grid, then the two tasks separately, then the levers, and finally the
list of things this benchmark cannot tell you.

\subsection{The design}

Three things vary, and nothing else:

\begin{itemize}
\item \textbf{Signal.} The 16 standardised abundances of \S\ref{sec:data};
  a 64- or 256-dimensional PCA projection of the APOGEE spectra
  (\S\ref{sec:spectral}); a 256-dimensional masked-autoencoder latent of the
  same spectra; the four kinematic observables
  (parallax, two proper motions, radial velocity); and concatenations of
  chemistry with kinematics.
\item \textbf{Method.} t-SNE, UMAP and EVoC, each followed by
  HDBSCAN$^*$ clustering where the method needs one, with the defaults of
  \textbf{Table~\ref{tab:config}}.
\item \textbf{Task.} \emph{Cluster-only separation}: cluster the known members
  among themselves and ask whether the clusters come out distinct, scored with
  \textbf{Equation~\ref{eq:merits}}. \emph{Field retrieval}: search the full
  sample for members and score with \textbf{Equation~\ref{eq:rp}}.
\end{itemize}

Three conventions keep the comparison honest, and all three are choices a
reader should check before believing any row:

\begin{enumerate}
\item \textbf{Same population.} Every arm in \textbf{Table~\ref{tab:honest}}
  is scored on the same 982 stars and the same 25 clusters. An earlier version
  of this table compared arms scored on 25 clusters against arms scored on 5,
  which is the mistake of \S\ref{sec:validation}.
\item \textbf{Seven seeds, reported as mean $\pm$ s.d.} The kinematics-only
  t-SNE row is deterministic, several other rows are not, and the difference
  is information rather than noise to be hidden.
\item \textbf{Two referees.} Field retrieval is scored against two independent
  definitions of ``member'': the kinematic selection of
  \S\ref{sec:data}, and the Simbad cluster memberships of the literature.
  Every conclusion below holds under both; the Simbad referee moves precision
  by $-0.01$ to $-0.08$ (the largest shifts are in the 64-dimensional PCA arms)
  and changes no ranking.
\end{enumerate}

\subsection{Task 1: separating the clusters from each other}

\textbf{Table~\ref{tab:honest}} is the central result of the workbook. Read it
by column first: the columns are methods, the rows are the signal the method
was given.

\begin{table}[htbp]
\centering
\small
\tabcolsep4pt
\begin{tabular}{@{}llccc@{}}
\toprule
signal & dim & t-SNE & UMAP & EVoC \\
\midrule
abundances                & 16  & $0.560 \pm 0.000$ & $0.578 \pm 0.011$ & $0.420 \pm 0.043$ \\
PCA of spectra            & 64  & $0.741 \pm 0.000$ & $0.754 \pm 0.007$ & $0.733 \pm 0.008$ \\
PCA of spectra            & 256 & $0.688 \pm 0.003$ & $0.685 \pm 0.014$ & $0.653 \pm 0.020$ \\
masked AE latent          & 256 & $0.740 \pm 0.002$ & $0.760 \pm 0.011$ & $0.714 \pm 0.036$ \\
kinematics only           & 4   & $0.942 \pm 0.000$ & $0.940 \pm 0.005$ & $0.876 \pm 0.012$ \\
\midrule
abundances, K-means ($K$ known) & 16 & \multicolumn{3}{c}{$0.449$ (single run)} \\
\bottomrule
\end{tabular}
\caption{Cluster-only homogeneity on the same 982 stars in 25 clusters, mean
$\pm$ standard deviation over seven seeds. Higher is better. The K-means row
is the classical baseline of \S\ref{sec:kmeans}, run on the 1\,002-row
member matrix with $K=25$ and $n_{\rm init}=10$ by
\texttt{article/scripts/kmeans\_baseline.py}: it is not an embedding, so it
has no column to sit in, and it is quoted for scale rather than for ranking.
The recovery fractions of the same arms, and the comparison with the published
result, are in \textbf{Table~\ref{tab:literature}}.}
\label{tab:honest}
\end{table}

Four things in that table are worth stating out loud.

\textbf{Chemistry is not nothing.} All three methods reach homogeneity well
above the trivial floor of $1/25 = 0.04$: the 16 abundances carry real cluster
information, and the embedding methods extract some of it. This is the
optimistic half of the chemical-tagging programme, reproduced on data that
post-date the paper that reported it.

\textbf{Chemistry is not enough.} The same methods on four kinematic
observables --- a quarter of the dimensions --- reach $0.94$. Nothing about
the abundance arms is competitive with kinematics, and no amount of embedding
repairs it: the four-dimensional signal is simply more informative about
which cluster a star belongs to.

\textbf{Dimension helps, up to a point.} Going from 16 abundances to the
64-dimensional PCA latent raises EVoC's homogeneity from $0.42$ to $0.73$ and
t-SNE's from $0.56$ to $0.74$. Going on to 256 principal components \emph{lower}
it --- $0.69$ for t-SNE, $0.65$ for EVoC --- which is the bias--variance trade
of \S\ref{sec:fundamentals} in its least forgiving form: the extra 192
directions carry variance, not signal, and every one of them is a direction in
which the field can become an apparent neighbour.

\textbf{One row of that table is not a measurement.} The abundances-only
t-SNE row is unstable under row order, as the audit of \S\ref{sec:validation}
showed: the same matrix gives $0.218$--$0.560$ depending on how the rows are
ordered, and the single-run baseline of \textbf{Table~\ref{tab:baseline}}
gives $0.292$ for what is nominally the same experiment. Any claim that rests
on the difference between $0.29$ and $0.56$ for t-SNE on raw abundances is
built on the input ordering.

\begin{table}[htbp]
\centering
\small
\tabcolsep5pt
\begin{tabular}{@{}lccc@{}}
\toprule
features & t-SNE & UMAP & EVoC \\
\midrule
abundances (16-d)          & 0.292 & 0.547 & 0.409 \\
abundances + kinematics    & \textbf{0.767} & 0.748 & 0.680 \\
\bottomrule
\end{tabular}
\caption{The single-run baseline of the DR19 re-run
(\texttt{cluster baseline}, no seed averaging, the full 1\,002-row member
matrix). It is included because it is the number our own documentation quotes,
and because comparing it with \textbf{Table~\ref{tab:honest}} is instructive:
concatenating kinematics lifts every method by $0.25$--$0.48$, and the
abundances-only t-SNE row differs between the two tables by more than the seed
spread of any other row.}
\label{tab:baseline}
\end{table}

\begin{table}[htbp]
\centering
\small
\tabcolsep5pt
\begin{tabular}{@{}llccc@{}}
\toprule
signal (excluding M~3) & t-SNE & UMAP & EVoC \\
\midrule
abundances (884 stars)              & $0.199$ (degenerate) & $0.581$ & $0.499$ \\
masked AE latent (884 stars)        & $0.711$ & $0.742$ & $0.658$ \\
\bottomrule
\end{tabular}
\caption{Ablation on the 982-star population: the globular cluster M~3
contributes $98$ of the stars ($884$ remain) --- or $154$ of the
1\,002 member rows of \textbf{Table~\ref{tab:clusters}} --- and it is the
largest globular in the sample, though not the largest cluster: M~67 supplies
$218$ stars and the Pleiades $113$. Remove M~3 and the abundances-only t-SNE arm does not
merely change, it collapses: two groups with $75\%$ of the stars in the
largest, zero stars labelled noise. The spectral latent is unaffected by the
same ablation. A result that depends on one cluster being present is not a
measurement of the method. The per-cluster counts of both populations are in
\texttt{results/population\_counts.csv}
(\texttt{python scripts/population\_counts.py --baseline-min-finite 9}).}
\label{tab:ablation}
\end{table}

\begin{figure}[htbp]\centering
  \fig[0.98]{benchmark_grid.png}
  \caption{The whole grid in one figure: method (rows) against signal
  (columns), for the two tasks. Read the left block for cluster-only
  separation and the right block for field retrieval, and notice that the
  two blocks do not agree on a winner --- no single row is best at both. This
  is the figure the lecture ends the benchmark on, and it is the reason
  \S\ref{sec:benchmark} refuses to name a best method.}
  \label{fig:benchgrid}
\end{figure}

\begin{figure}[htbp]\centering
  \fig[0.86]{headtohead_pca.png}
  \caption{The head-to-head on $64$-dimensional PCA spectra, cluster-only
  separation, seven seeds. The error bars are the seed spread: the masked-AE
  and PCA-64 arms overlap, which is the tie reported in
  \S\ref{sec:spectral} and the reason the spectral result is stated as
  ``spectra beat abundances, and do not beat PCA''.}
  \label{fig:headtohead}
\end{figure}

\subsection{Task 2: finding members in the field}

Now the operational question. Take 30\,107 stars (the 25 clusters plus a
stratified sample of the field), cluster them, and ask how many of the known
members were recovered and how much of what was returned was real.

\begin{table}[htbp]
\centering
\small
\tabcolsep4pt
\begin{tabular}{@{}llccc@{}}
\toprule
signal & method & recall & precision & kNN purity \\
\midrule
abundances (16-d) & t-SNE & 0.201 & 0.240 & 0.124 \\
abundances (16-d) & UMAP  & 1.000 & 0.001 & 0.075 \\
abundances (16-d) & EVoC  & 0.557 & 0.003 & --- \\
\midrule
PCA 64-d & t-SNE & 0.417 & 0.631 & 0.528 \\
PCA 64-d & UMAP  & 0.374 & 0.592 & 0.434 \\
PCA 64-d & EVoC  & 0.585 & 0.003 & --- \\
\midrule
masked AE 256-d & t-SNE & 0.418 & 0.619 & 0.504 \\
masked AE 256-d & UMAP  & 0.385 & 0.513 & 0.419 \\
masked AE 256-d & EVoC  & 0.255 & 0.005 & --- \\
\bottomrule
\end{tabular}
\caption{Field retrieval, all-sky, on the same 30\,107 stars and the same 25
clusters, scored against the kinematic membership (978 members). Recall and
precision are the best-overlap definition of \textbf{Equation~\ref{eq:rp}};
kNN purity is the parameter-free cohesion score of \S\ref{sec:knn}, and it is
undefined for EVoC because that arm returns no embedding to measure it in.
Scoring against the Simbad referee instead shifts every precision by
$-0.01$ to $-0.08$ and changes no ranking.}
\label{tab:field}
\end{table}

The three lessons here organise most of what the rest of this workbook says
about chemical tagging in practice.

\textbf{Recall without precision is a statement about the field, not about
the clusters.} UMAP on raw abundances recovers \emph{every} true member ---
at a precision of $0.001$. It has not found the clusters; it has labelled a
large part of the field as cluster, and the best-overlap protocol dutifully
reports that the members are inside the resulting group. EVoC does the same
thing more politely ($0.557$ recall, $0.003$ precision). This is the trap that
the literature's habit of quoting a single number walks into, and it is why
\textbf{Table~\ref{tab:field}} prints three.

\textbf{Spectra read directly beat element ratios, and it took a chemical
ablation to be sure.} The 64-dimensional PCA latent and the 256-dimensional
masked-AE latent both reach recall $\approx 0.42$ at precision
$\approx 0.62$ from the abundances' $0.20$/$0.24$. Two independent
representations of the same spectra agree with each other to within
$0.002$ in recall and $0.012$ in precision --- which is the honest way to say
that the improvement is real while the difference between the $64$- and
$256$-dimensional spectral rows is not.

\textbf{The gap to the ceiling is the story.} Kinematics alone recover $0.93$
of the members at a precision limited by velocity doppelg\"angers rather than
by the method (\S\ref{sec:fundamentals}). Chemistry, read as well as we can
read it today, gets to $0.42$ recall. Strong chemical tagging --- recovering a
dispersed cluster with no kinematic help --- is not what these numbers
demonstrate; what they demonstrate is that modern self-supervised spectral
representations find \emph{something} real about a cluster in a spectrum, and
that the something is diluted by a factor of two relative to parallaxes and
proper motions.

\subsection{The curse of dimensionality, measured}

Concatenating a high-dimensional, weakly informative signal onto a
low-dimensional, strongly informative one does not add their information. It
dilutes the strong signal, because distance is shared across dimensions. The
DR17 study measured this on 878 members in region mode with the kinematic
labels held out of the ground truth:

\begin{table}[htbp]
\centering
\small
\tabcolsep5pt
\begin{tabular}{@{}lccc@{}}
\toprule
signal (dims) & cluster homogeneity & field recall & field precision \\
\midrule
abundances (16)          & 0.25 / 0.51 / 0.51 & 0.73 & 0.47 \\
spectral latent (256)    & 0.40 / 0.69 / 0.64 & 0.72 & 0.51 \\
+ kinematics (20 / 260)  & 0.62 / 0.74 / 0.72 & 0.83 / 0.80 & 0.63 / 0.61 \\
kinematics only (4)      & \textbf{0.97 / 0.96 / 0.92} & \textbf{0.93} & 0.63--0.67 \\
\bottomrule
\end{tabular}
\caption{Signal comparison on the DR17 data (three numbers per cell are the
t-SNE / UMAP / EVoC homogeneity). Kinematics alone near-solve the cluster-only
task and set the field-recall ceiling; adding chemistry on top improves the
cluster-only score slightly and never improves recall beyond what kinematics
already achieved. The DR19 re-run reproduces the pattern with different
absolute values (kinematics $0.94$, abundances $0.56$,
abundances$\,+\,$kinematics $0.77$ over the same 982 stars).}
\label{tab:curse}
\end{table}

The mechanism is worth stating precisely, because it is the reason this
workbook spends three chapters on geometry rather than one on algorithms. A
distance in a 20-dimensional space is an average over 20 coordinate
differences. If 16 of those coordinates are chemistry and carry weak
information about cluster identity, while 4 are kinematics and carry strong
information, then the strong coordinates are outvoted$\,$--- not mis-weighted,
outvoted, because no reweighting recovers information that the metric has
already blurred. Rescaling chemistry \emph{up} to compensate costs precision
in the field, as \S\ref{sec:benchmark}'s lever table shows. The correct
response is architectural, not numerical: cluster on the strong signal, then
apply the weak signal as a rejection test. That is the two-stage pipeline, and
it is the single most successful idea in this project.

\subsection{Levers, and a lesson about bug fixes}

The DR17 study swept the pipeline's knobs on one cluster (M~67) inside a
$30^\circ$ region, with $5\,000$ stars and 183 referee members. Every run
logged its full configuration. The table is worth reading as a list of
\emph{hypotheses}, not results, and the section that follows explains why.

\begin{table}[htbp]
\centering
\small
\tabcolsep4pt
\begin{tabular}{@{}lcccc@{}}
\toprule
lever & recall & prec. & UMAP recall & verdict \\
\midrule
baseline                                   & 0.087 & 0.117 & 0.415 & --- \\
\texttt{min\_cluster\_size} $5\!\to\!10$   & 0.404 & 0.129 & 0.525 & best single \\
perplexity $30\!\to\!50$                   & 0.383 & 0.124 & 0.415 & free win \\
perplexity $30\!\to\!10$                   & 0.033 & 0.353 & 0.415 & precision for recall \\
\texttt{SNR\_MIN} $100\!\to\!150$          & 0.296 & 0.089 & 0.310 & recall up, precision down \\
dwarf-only sample                          & 0.149 & 0.165 & 0.158 & best purity \\
drop S, K, V, Mn                           & 0.055 & 0.256 & 0.038 & hurts recall \\
UMAP \texttt{metric=cosine}                & 0.087 & 0.117 & 0.044 & hurts UMAP \\
UMAP \texttt{n\_neighbors} $=5$            & 0.087 & 0.117 & 0.071 & worse \\
UMAP \texttt{n\_neighbors} $=30$           & 0.087 & 0.117 & 0.410 & $=$ baseline \\
element weights (after a bug fix)          & \textbf{0.437} & \textbf{0.158} & 0.339 & best lever \\
weights $+$ size floor                     & 0.432 & 0.160 & 0.525 & not additive \\
\bottomrule
\end{tabular}
\caption{DR17 lever sweep (M~67, $30^\circ$ region, $5\,000$ stars, 183
referee members, seed 42). The recall and precision columns are t-SNE; the
last two rows exist because of a bug: the element-weighting switch was
originally a no-op, since the weights were applied and then undone by the
standardisation step that followed. Reordering the two operations turned the
largest apparent non-effect into the largest effect in the table. Two rules of
thumb survive from this sweep --- the size floor of $5$ fragments the field,
and a perplexity below $N/100$ starves the embedding --- and one warning: the
best levers do not compose.}
\label{tab:levers}
\end{table}

The lesson is not that hyperparameters matter, which is known, but that the
\emph{largest single effect in a sweep can be a bug}. A configuration flag that
does nothing is worse than a missing flag, because it produces a negative
result that looks like an experiment. The DR17 re-check of the same lever on
DR19 data did not reproduce the recall gain, which is exactly what the protocol
of \S\ref{sec:validation} predicts for a lever that was selected on one
cluster at one seed: selecting on a metric and reporting the same metric is
not validation. The DR19 region sweep of the same switch confirms the warning
from the other direction --- a recall that climbs to $0.96$ while precision
falls to $0.05$ is a cluster that has grown to swallow the field.

\subsection{Where this sits in the literature: Casamiquela et al. (2021)}
\label{sec:literature}

Every number so far in this chapter has been scored against our own baselines.
The closest published measurement is \citet{Casamiquela:21}, and it is close
enough to sit in the same table: the same metric triple, the same task, and a
sample built to be the best case that spectroscopy can produce. They measured 16
elements differentially in 175 red-clump members of 31 thin-disc open clusters,
with an internal coherence of about $0.03$ dex per element and \emph{no field
stars}. They then clustered the abundances directly with HDBSCAN --- no
embedding, no projection --- and fine-tuned its free parameters to maximise the
number of correctly recovered clusters (their Sect. 4.2). They call the sample
itself their \emph{best-case scenario}. Their published result is $h = 0.49$, $c = 0.63$, $V = 0.55$,
$\mathrm{RF}_{40} = 0.29$ (9 of 31 clusters) and $\mathrm{RF}_{70} = 0.03$ (one
cluster, NGC~2682); their own summary of that run is that more than 70\,\% of the
groups it finds are statistical, mixing stars from several real clusters.

\begin{table}[htbp]
\centering
\small
\tabcolsep5pt
\begin{tabular}{@{}lcccccc@{}}
\toprule
sample & $n_\star$ / $n_{\rm cl}$ & $h$ & $c$ & $V$ & $\mathrm{RF}_{40}$ & $\mathrm{RF}_{70}$ \\
\midrule
\citet{Casamiquela:21}, published                  & 175 / 31 & 0.49  & 0.63  & 0.55  & 0.290 & 0.030 \\
their clustering step, our stars (raw dex)         & 982 / 25 & 0.499 & 0.416 & 0.454 & 0.080 & 0.040 \\
their clustering step, our stars (standardised)    & 982 / 25 & 0.475 & 0.407 & 0.438 & 0.000 & 0.000 \\
\midrule
our pipeline: t-SNE                                & 982 / 25 & 0.560 & 0.601 & 0.580 & 0.320 & 0.080 \\
our pipeline: UMAP                                 & 982 / 25 & 0.578 & 0.572 & 0.575 & 0.326 & 0.006 \\
our pipeline: EVoC                                 & 982 / 25 & 0.420 & 0.529 & 0.466 & 0.189 & 0.006 \\
our pipeline: K-means ($K$ known)                  & 982 / 25 & 0.454 & 0.490 & 0.471 & 0.080 & 0.000 \\
\midrule
open clusters only, UMAP                           & 665 / 18 & 0.573 & 0.422 & 0.486 & 0.405 & 0.040 \\
open clusters only, t-SNE                          & 665 / 18 & 0.471 & 0.458 & 0.465 & 0.222 & 0.000 \\
open clusters only, EVoC                           & 665 / 18 & 0.520 & 0.419 & 0.464 & 0.214 & 0.048 \\
\bottomrule
\end{tabular}
\caption{The published comparison. Rows 2--3 are \textit{their} clustering step
--- HDBSCAN on the abundances, \texttt{min\_cluster\_size}$=2$, their
$4 \times 7$ parameter grid, best grid point by $\mathrm{RF}_{40}$ --- run on
our stars; their configuration is defined in dex, so the raw-abundance row is
the faithful replication and the standardised row is a rescaling check. Their
$4 \times 7$ grid never reaches $\mathrm{RF}_{40} > 0.11$ on our stars. Rows
4--10 are our arms, mean over seven seeds (the $h$, $c$, $V$ columns repeat
\textbf{Table~\ref{tab:honest}}); the open-cluster subset is the closest
analogue of their sample. Chance levels, from shuffling the true labels into
groups of the observed sizes: for their sample shape $V = 0.51$ ($0.44$--$0.51$
across the three partition shapes tested) and $\mathrm{RF}_{40} = 0.04$ (95th
percentile $0.10$); for their clustering step run on \emph{our} stars $V = 0.29$
--- their rule cuts 165 groups, and many small groups inflate the null --- with
$\mathrm{RF}_{40} \le 0.001$; for our own arms $V = 0.06$--$0.13$ and
$\mathrm{RF}_{40} \le 0.001$. Source:
\texttt{results/casamiquela\_comparison*.csv} and
\texttt{results/casamiquela\_protocol.csv}; the second one also carries the grid
search.}
\label{tab:literature}
\end{table}

Four things follow from that table.

\textbf{Their clustering step recovers less on our stars than on theirs.} Run
their configuration on our 982 stars and the recovery fraction falls to
$0.08$ ($0.11$ on the open clusters alone, both being Berkeley~17 and
NGC~1245), against $0.29$ on their own sample; under the standardised variant it
returns 168 groups and recovers nothing at all. Two differences remain that we
cannot separate: our abundances are survey-quality --- $\sim\!0.05$ dex per
element against their $0.03$ --- and our clusters are not a single evolutionary
stage.

\textbf{Our two-stage pipeline lands in the same range as their one-stage
pipeline.} t-SNE and UMAP on 16 survey abundances, with default HDBSCAN, recover
$0.32$ and $0.33$ of the 25 clusters, and on the 18 open clusters alone --- the
subset closest to their sample --- UMAP reaches $0.41$. The honest reading is not
that we beat them. It is that embed-then-cluster on coarser abundances reaches
the same recovery as cluster-the-abundances on the best abundances available,
and that neither gets past about a third of the clusters.

\textbf{The failures have the same shape.} Their surviving clusters sit at the
edges of the abundance distribution --- NGC~6705, NGC~2682, and NGC~2420, which
they call the most metal-poor cluster in their sample. NGC~2420 is a regular
recovery of ours too: UMAP recovers it in every seed of the open-cluster subset
and EVoC in six of seven. Our arms otherwise recover the same handful of
clusters in every seed --- Berkeley~17, IC~166, M~71, NGC~1245, NGC~6819,
NGC~7789, the Pleiades --- and that set contains M~3, which \S\ref{sec:clusters}
introduces as the metal-poor outlier of the sample. Whether this is the same
edge-of-distribution mechanism their paper describes is a hypothesis this
workbook does not settle: member evolutionary state is a confound we have not
yet removed, which is what the \texttt{DWARF\_ONLY} lever of
\S\ref{sec:assignment} is for.

\textbf{Their field-star test is the published analogue of our Task 2, and it
agrees.} Adding field stars to their sample leaves one cluster above the 40\,\%
threshold (NGC~188) and none at 70\,\%. Our field retrieval
(\textbf{Table~\ref{tab:field}}) agrees in direction and goes further: the best
abundance-only arm reaches recall $0.20$ at precision $0.24$, and the arms that
reach high recall do it by labelling most of the field as cluster. Their element
tests agree with ours too: restricting the set to 8 or 7 elements did not improve
their recovery, and our dropped-element row in \textbf{Table~\ref{tab:levers}}
loses recall as well.

\subsection{What this benchmark cannot tell you}
\label{sec:limitations}

Five limits, in descending order of how much they should trouble a reader.

\begin{enumerate}
\item \textbf{The sample is small and uneven.} 982 stars across 25 clusters is
  a few tens of stars per object, and the distribution is wide: M~67 supplies
  218 stars and the Pleiades 113, while five clusters supply fewer than ten each
  once the duplicate rows are collapsed (Berkeley~17, NGC~1798, King~5,
  NGC~2158, King~7 --- \textbf{Table~\ref{tab:clusters}} counts rows, and lists
  three). \textbf{Table~\ref{tab:ablation}}
  shows what a single well-populated cluster can do to a result, and every macro
  average in this chapter is therefore an average over a handful of clusters
  that are not equally measured. The member frame is itself
  configuration-dependent: the 1\,002 rows counted in
  \textbf{Table~\ref{tab:clusters}} need a finite-abundance floor of nine out
  of sixteen elements, while the shipped default of eight keeps 1\,215 rows ---
  so the saved baseline scores are not reproducible from a bare
  \texttt{cluster baseline} run
  (\texttt{python scripts/population\_counts.py --baseline-min-finite 9}).
\item \textbf{The member matrix contains duplicated stars.} 298 rows share an
  \texttt{APOGEE\_ID} with another row, the duplicates differ by
  $\sim\!0.04$ dex per element, and they inflate the largest clusters most.
  Until the sample is collapsed to unique stars, every macro average in this
  chapter is a weighted average over something that is not quite a
  star.
\item \textbf{Some rows are not measurements.} The abundances-only t-SNE row
  is row-order unstable (\S\ref{sec:validation}); the ablation row is
  degenerate by construction. Both are printed rather than dropped, because
  the pattern of which arms are fragile is itself a result: the
  spectral-latent arms are stable, and the raw-abundance t-SNE arm is not.
\item \textbf{The external comparison is one pipeline, not a survey of the
  literature.} \S\ref{sec:literature} runs \citet{Casamiquela:21}'s clustering
  step on our stars and their metric triple on our scores, which is as far as a
  like-for-like comparison goes without their star list. We do not re-run the
  graph-attention autoencoder of \citet{Spina:25}, which works on
  $\sim\!47\,000$ APOGEE stars and deliberately adds kinematics and age to the
  chemistry, so any claim of state-of-the-art in this space should be read as
  unaudited.
\item \textbf{The counter-result stands, and it is not the same measurement.}
  \citet{Hogg:16} recovered phase-space structure by clustering 15 abundances
  with K-means, and our K-means row in \textbf{Table~\ref{tab:honest}} matches
  EVoC's homogeneity. On the recovery fraction it is the weakest arm of the
  table ($0.08$ against t-SNE's $0.32$), which is the difference between a tidy
  partition and a recovered cluster: $K$-means is told how many groups to make,
  so its homogeneity is high by construction, while the density methods have to
  find the groups. Whatever this workbook demonstrates about the limits of
  abundance-space clustering, it is a statement about APOGEE's ASPCAP
  precision, these 25 clusters, and these methods --- not a statement that
  chemical tagging does not work.
\end{enumerate}

\begin{issues}[EXERCISES]
\begin{enumerate}
\item Reproduce one row of \textbf{Table~\ref{tab:honest}}: pick a signal and
  a method, run it over the seven seeds, and report mean $\pm$ standard
  deviation. Compare with the published row. If your numbers differ, is the
  cause the population, the seed set, or a default?
\item \textbf{Table~\ref{tab:field}} reports precision for the kinematic
  referee. Recompute it against the Simbad memberships and plot the
  difference per cluster. Which clusters are most referee-sensitive, and does
  that correlate with their age or their [Fe/H]?
\item The ablation of \textbf{Table~\ref{tab:ablation}} removes M~3. Repeat
  it for the other clusters with more than 100 rows (M~67, NGC~6819, NGC~2243)
  and for the duplicates-cleaned matrix. How much of the instability survives
  cleaning?
\item \S\ref{sec:literature} runs their pipeline on our stars but not ours on
  theirs. Their sample reachable star list is not public in a form we can
  re-run, so state what you would need from the other authors --- and what a
  re-run on their stars would change about the conclusion of
  \textbf{Table~\ref{tab:literature}}.
\end{enumerate}
\end{issues}

% 14 · Spectral representations — reading the whole spectrum
\section{SPECTRA INSTEAD OF ELEMENT RATIOS}
\label{sec:spectral}

\deckslides{slide 47, the masked-autoencoder block and its four steps.}

Everything in \S\ref{sec:benchmark} treats a star as a 16-dimensional vector of
abundances. That vector is a lossy summary of the observation it came from: a
single APOGEE spectrum has 8\,575 wavelength pixels, and the abundance
pipeline that turns those pixels into 16 numbers assumes a one-dimensional
model atmosphere, a set of atomic and molecular line data, and a fitting
procedure that was tuned on the Sun. This chapter asks the obvious question ---
if the spectra are the data, why cluster the summary? --- and reports what
happens when we do.

\begin{marginnote}[]
\entry{Masked autoencoder}{A self-supervised model that hides part of its
input, encodes the visible part, and is trained to reconstruct the hidden
part. It learns the structure of the data without any labels. In 1-D spectra,
\emph{how} you mask decides what it can learn.}
\entry{Self-supervised}{Trained with a loss defined by the data itself
(reconstruction) rather than by external labels. The property we care about
here is negative: a self-supervised latent never saw an element ratio or a
cluster membership, so it cannot be circular.}
\end{marginnote}

\subsection{The baseline: PCA of the spectra}

The honest first experiment is principal component analysis. \citet{GarciaDias:19}
applied eight clustering algorithms to the 13 APOGEE abundances of each star
under four dimensionality-reduction strategies, and this chapter asks the same
question one level down: which representation of the observation keeps the
cluster information. Reading spectra with
learned models has a longer history in this field --- the Cannon builds a
data-driven spectral model from labels \citep{Ness:15}, and convolutional
networks do the same job without the labels' linearity assumption
\citep{Ting:19} --- and the question this chapter asks is which of those
representations keeps the cluster information. Keep 64 components and
the classification results of \S\ref{sec:benchmark} already improve on the
abundances; keep 256 and they get slightly worse. PCA is the right baseline and
the wrong tool for one reason that is worth internalising: it maximises
\emph{variance}, and in a stellar spectrum the variance is dominated by things
that have nothing to do with chemistry --- the temperature-sensitive continuum
shape, gravity-sensitive line profiles, interstellar reddening, and the
signal-to-noise of the observation. A principal component that encodes
$T_{\rm eff}$ is useful for almost every other problem in stellar
spectroscopy and actively unhelpful here, because a hot field star and a hot
cluster member will share it.

\subsection{The model: a masked spectral autoencoder}

Our proposal \citep{He:22} is to train a self-supervised latent that has to
predict the spectrum's hidden parts:

\begin{enumerate}
\item Take the continuum-normalised spectrum, 8\,575 pixels.
\item Mask 50\% of the wavelength axis, in \emph{contiguous blocks} of 200
  pixels --- roughly a spectral line region at APOGEE's resolution.
\item Encode the visible pixels with five convolutional blocks, each halving
  the wavelength axis, with channel widths from 1\,024 down to 64, ending in a
  global average pool and a 256-dimensional linear bottleneck.
\item Reconstruct the masked pixels with a mirrored decoder --- and compute the
  loss \emph{only where the input was hidden}.
\item Discard the decoder. Keep the encoder's 256-dimensional output and feed
  it to the clustering pipeline exactly as the abundances were.
\end{enumerate}

The model has 8.6 million parameters and trains in under an hour on one GPU on
39\,945 field spectra plus 994 member spectra, with early stopping. Nothing in
the loss function mentions abundances, clusters, kinematics, or age.

\begin{figure}[htbp]\centering
  \fig[0.60]{mae_arch.png}
  \caption{The architecture: five strided convolutional blocks compress the
  8\,575-pixel spectrum by a factor of 32 while widening the channel count, a
  global average pool removes the wavelength axis, and a linear layer produces
  the 256-dimensional latent. The decoder is a mirror of the encoder and is
  discarded at inference.}
  \label{fig:maearch}
\end{figure}

\subsection{Why the masking is contiguous}

This is the one design decision in the chapter that a reader should be able to
justify from first principles, and it is where a naive implementation fails.

Mask single pixels at random and a convolutional network does not have to
learn chemistry: the masked pixel's immediate neighbours are visible, and the
spectrum is smooth on the scale of a few pixels, so copying from the left and
right solves the task. The reconstruction loss falls, the latent looks
respectable, and the model has learned to interpolate. Mask \emph{contiguous
200-pixel blocks} and the visible pixels adjacent to the hole are far enough
away, in wavelength terms, to be in different lines; the model must use the
global shape of the spectrum --- which line regions are present, and how deep
--- to fill in the gap. That is where element ratios live.

\begin{figure}[htbp]\centering
  \begin{tabular}{@{}c@{\hskip6pt}c@{}}
    \fig[0.45]{mae_why_blocks.png} &
    \fig[0.45]{mae_step1_mask.png}
  \end{tabular}
  \caption{Left: why the masking is contiguous --- a single-pixel hole is
  filled by interpolation from its immediate neighbours, while a 200-pixel
  block removes a whole line region and forces the model to use context.
  Right: the actual mask used (top) and the residual the model must predict
  (bottom), with the masked region shaded.}
  \label{fig:maemask}
\end{figure}

\begin{figure}[htbp]\centering
  \begin{tabular}{@{}c@{\hskip6pt}c@{}}
    \fig[0.45]{mae_step2_encode.png} &
    \fig[0.45]{mae_step3_recon.png}
  \end{tabular}
  \caption{Encoding and reconstruction. Left: only the visible 50\% of the
  spectrum enters the encoder. Right: the decoder's reconstruction of the
  hidden half, against the truth, with the loss computed only on the hidden
  pixels. The model is never rewarded for the parts it can see --- which is
  what makes its latent a description of the spectrum rather than a copy of
  it.}
  \label{fig:maeenc}
\end{figure}

\begin{figure}[htbp]\centering
  \fig[0.55]{mae_step4_latent.png}
  \caption{The 256-dimensional latent, projected. The spectral latent carries
  the abundance structure of the members because it had to reconstruct the
  lines those abundances are measured from --- not because anyone told it she
  was looking at C, N, O or Fe.}
  \label{fig:maelatent}
\end{figure}

\subsection{What it achieves}

On cluster-only separation it is the strongest signal we have
(\textbf{Table~\ref{tab:honest}}): homogeneity $0.740 \pm 0.002$ (t-SNE),
$0.760 \pm 0.011$ (UMAP), $0.714 \pm 0.036$ (EVoC) on the 982-star population,
against abundances $0.560$/$0.578$/$0.420$. In the single-run baseline it
reaches $0.789$/$0.874$/$0.770$ on the full 1\,002-row matrix, and
concatenating it with the abundances gives $0.879$ --- the highest cluster-only
score measured anywhere in this project. It also has the most important
non-empirical property: because it never saw an abundance label, it cannot be
circular in the way a supervised latent is.

On field retrieval it is a different story, and the difference is the most
useful result in this chapter:

\begin{itemize}
\item \textbf{Facing the field, the masked latent does not beat PCA-64.} In
  \textbf{Table~\ref{tab:field}} the two are
  $0.418$/$0.619$ and $0.417$/$0.631$: a tie to within the third decimal.
  Two very different representations of the same spectra --- one unsupervised
  and linear, one self-supervised and deep --- land on the same operating
  point.
\item \textbf{The tie is the finding.} If a 64-dimensional linear projection
  reaches the same field-retrieval precision as a 8.6-million-parameter
  convolutional latent, then what is limiting the precision is not the
  representation capacity of the model. It is the information in the spectra
  themselves, or the way our metric uses it.
\item \textbf{EVoC disagrees, and we report both.} On the same cluster-only
  task with EVoC as the clusterer, the supervised CNN reaches
  $0.80 \pm 0.03$ on the four-cluster sample against the masked latent's
  $0.70 \pm 0.08$ --- the supervised model wins, and the two error bars
  overlap. We described the self-supervised result as a
  t-SNE/UMAP result, not as a universal one, and withdrew the stronger
  sentence.
\end{itemize}

\begin{figure}[htbp]\centering
  \fig[0.72]{paired_control.png}
  \caption{The paired control that ended the ``two times purer'' claim, which
  is the reason this figure is in the workbook. Each star is embedded through
  both pipelines; the panel shows the distance from a star to \emph{itself}
  across data products against the distance to a different random star within
  one product. The mean is indicated. When a star is further from its own
  second embedding than from a stranger, the latent is encoding something
  about the data product rather than about the star.}
  \label{fig:paired}
\end{figure}

\subsection{The error we made, in full}
\label{sec:provenance}

A masked autoencoder trained on a mixed set of spectra will learn whatever
distinguishes them, and our first training set was mixed in a way we did not
notice. The field sample was 100\% raw \texttt{apStar} spectra, while 91\% of
the member sample was continuum-normalised \texttt{aspcapStar} output, with
median flux levels differing by a factor of $\sim\!6\,000$. Along that boundary
the labels themselves are the difference.

The consequences, all measured:

\begin{itemize}
\item A linear probe recovers the data product from the latent at
  \textbf{AUC 0.999} on members alone, using 223 of the 256 dimensions
  ($p < 10^{-6}$).
\item A one-bit flag (``which product is this?'') separates members from field
  stars at \textbf{AUC 0.957} --- essentially the whole latent score, with no
  chemistry involved.
\item The paired control of \textbf{Figure~\ref{fig:paired}} is decisive: for
  253 stars embedded through both pipelines, a star sits $1.70\times$ further
  from its own other-product embedding than from a random different star
  within one product (cosine distance $0.389$, with $75\%$ of the offset in
  one shared direction --- the signature of a single global nuisance axis).
\end{itemize}

We first diagnosed this as a data-release effect (DR17 members, DR19 field).
That was wrong, and the check that disproved it is worth repeating: DR19
reanalyses and includes DR17 --- 717\,689 rows carry
\texttt{release=`dr17'} --- and DR19 serves a uniform \texttt{mwmStar}
spectrum for \textbf{738 of 738} members in the backfilled list. The mismatch
was between two products on the same sky, at the same epoch, from the same
spectrograph.

So the field-retrieval row for the mixed latent is not interpretable, and the
``twice as pure'' precision claim that depended on it is withdrawn. What
survives is narrower and solid: the cluster-only homogeneity, which is computed
between members and therefore within one product, holds on a single-product
subsample ($0.612$/$0.686$/$0.575$ against the abundances' $0.484$/$0.550$/
$0.488$ on the same stars). The fix is not a caveat but a re-run: download
uniform \texttt{mwmStar} spectra for every member and field star, re-embed, and
repeat the benchmark. That run is staged in the repository --- 736
\texttt{mwmStar} spectra are already downloaded and waiting --- and its results
belong in the next version of this workbook rather than in this one.

\begin{issues}[EXERCISES]
\begin{enumerate}
\item Train the same autoencoder twice, once with single-pixel random masking
  and once with contiguous 200-pixel blocks, and compare the reconstruction
  loss on a held-out set. Then compare the two latents on the cluster-only task.
  Which model wins on reconstruction, and which on chemistry? Explain the
  disagreement.
\item Reproduce the linear probe of \S\ref{sec:provenance}: fit a logistic
  regression from a 256-dimensional latent to a binary flag identifying the
  data product, and report the AUC on held-out member stars. Now repeat it for
  a latent trained on a single product. What is the AUC you would consider
  acceptable before publishing a field-retrieval number?
\item The paired control uses 253 stars with embeddings through both pipelines.
  Design an equivalent control for a study that only has one product, and
  state what it can and cannot detect.
\item The masked latent and PCA-64 tie on field retrieval. Working from the
  two latents, characterise what each one encodes: fit a linear model from
  each latent to $T_{\rm eff}$, $\log g$ and [Fe/H] in turn, and compare the
  cross-validated errors. Which representation spends more of its capacity on
  chemistry?
\end{enumerate}
\end{issues}

% 15 · Physics — what membership is actually for
\section{WHAT MEMBERSHIP IS FOR}
\label{sec:physics}

\deckslides{slide 47, the isochrone and two-survey blocks.}

Every number so far has been a score against labels. This chapter closes the
loop: membership lists are not the product, they are the input to a measurement
of what the cluster \emph{is} --- its age, its distance, its chemistry --- and
those measurements are what chemical tagging has to deliver if it is to matter
for Galactic archaeology \citep{Freeman:02, BlandHawthorn:16}.

\begin{marginnote}[]
\entry{Isochrone}{A theoretical curve in the colour--magnitude diagram giving
the luminosity and colour of every mass of a single-age, single-metallicity
population. Fitting it to a cluster's stars returns age, distance and
reddening.}
\entry{Distance modulus}{$m - M = 5\log_{10} d - 5$, the difference between
apparent and absolute magnitude. It is the fitted distance in a CMD; a
0.1 mag error is a 5\% error in distance.}
\end{marginnote}

\subsection{The measurement problem}
\label{sec:measurement}

\begin{figure}[htbp]\centering
  \begin{tabular}{@{}c@{\hskip6pt}c@{}}
    \fig[0.44]{sky_ngc188.png} & \fig[0.44]{sky_ic166.png}
  \end{tabular}
  \caption{The two ends of the sample, which is why the same pipeline behaves
  differently on different clusters. Left: NGC~188 --- old ($7.6$ Gyr), well
  populated, chemically distinct from the field. Right: IC~166 --- sparse and
  contaminated, with several member rows whose \texttt{APOGEE\_ID} is blank
  (\S\ref{sec:validation}). In the first, the isochrone fit has hundreds of
  stars on a clean sequence; in the second, every contaminant moves the fit.}
  \label{fig:skies}
\end{figure}

Fitting an isochrone grid to a cluster means searching a three-dimensional
parameter space (age, metallicity, distance modulus) for the combination that
best reproduces the observed colour--magnitude diagram, with the reddening as a
fourth, partly degenerate axis. We do it with the PARSEC grids through
\texttt{ASteCA} \citep{Perren:15}, sampling the posterior with \texttt{emcee}
\citep{ForemanMackey:13}. Two properties of that fit dominate everything that
follows.

\textbf{Age, metallicity and reddening are partially degenerate.} In the
regions of the CMD our clusters occupy, an older, more metal-rich population
at larger reddening can mimic a younger, metal-poor one at smaller reddening.
The degeneracy is a property of the isochrone geometry, not of the data or the
sampler: the posterior stays broad in the degenerate direction even when every
star in the cluster is a perfect member. In the two-survey fit below, the
age posterior width is $\sigma_{\log a} \approx 0.97$ --- nearly a factor
of ten in age --- for a cluster whose distance is pinned to 0.01 mag.

\textbf{Contaminants bias the parameters in a specific direction.} Field stars
in a cluster's colour--magnitude diagram occupy the same locus as the cluster's
own stars in some mass ranges and a different one in others, and the bias they
introduce is not random. Adding field stars to a metal-poor globular cluster
lowers the inferred age; adding them to a young cluster raises it. Membership
quality is therefore a parameter-estimation problem, not a bookkeeping one.

\subsection{Distances from red clump giants}

The most robust single measurement available to us is the distance, because
red clump giants have nearly constant absolute magnitude: the median $K$-band
magnitude of the members that fall in the $J-K$ clump slice is a standard
candle with a known calibration. For the clusters where our membership is
deep enough to have red clump members, the distances agree with the
literature:

\begin{table}[htbp]
\centering
\small
\tabcolsep6pt
\begin{tabular}{@{}lccc@{}}
\toprule
cluster & $dm$ (clump) & $dm$ (literature) & residual \\
\midrule
NGC~6819 & 11.92 & 11.90 & $+0.01$ \\
NGC~2243 & 13.08 & 13.25 & $-0.16$ \\
NGC~7789 & 11.68 & 11.27 & $+0.41$ \\
\bottomrule
\end{tabular}
\caption{Red-clump distance moduli from member giants, against literature
values. Two of the three agree to better than 0.2 mag, which for a membership
list of a few dozen giants assembled from a survey without any prior
membership information is the strongest validation in this chapter: the
members we find are located where the clusters are.}
\label{tab:dm}
\end{table}

\subsection{The case that decided us: NGC~2243 from two surveys}

NGC~2243 is the demonstration that chemical tagging is a two-survey problem.
APOGEE, being a near-infrared survey of the Galactic disc, observes a
cluster's \emph{giants}: bright, plentiful, and excellent for distances and
chemistry. GALAH observes at visible wavelengths and reaches the
\emph{main sequence and turnoff}, which is where age sensitivity lives. Neither
survey alone constrains both.

The sequence we ran is the one a workshop should reproduce:

\begin{enumerate}
\item Take the APOGEE member giants and fit an isochrone. The fit returns an
  age of 1.59 Gyr against a literature value of 1.07 Gyr --- wrong by half ---
  with a broad distance posterior ($\sigma_{dm} = 0.111$).
\item Use the red-clump distance (13.08 mag, agreeing with the literature
  value of 13.25 to within 0.17) as a tight prior, and select GALAH stars
  kinematically around the resulting centroid. At 4.7 kpc the parallax error
  rivals the parallax itself, so a relaxed parallax cut is required --- a
  detail that took a re-run to get right.
\item Fit the combined main-sequence-plus-giants sample: age $0.98$ Gyr
  (residual $0.09$), $\sigma_{dm} = 0.086$.
\end{enumerate}

A factor of six improvement in the age residual, from the same cluster and the
same survey data, by adding the stars that are sensitive to age and prior-ing
the distance with a measurement that is not. The age posterior remains broad
($\sigma_{\log a} \approx 0.97$), which is the degeneracy of \S 15.1 making
itself felt --- but the mode is right, and the distance is a fifth as
uncertain.

\subsection{The finale: better membership, better parameters}

The clearest scientific payoff of the whole project is a single cluster. Fit
an isochrone to three different membership lists for M~67 --- the literature
values are age $2.82$ Gyr, $dm = 9.54$, $[{\rm Fe/H}] = 0.03$:

\begin{table}[htbp]
\centering
\small
\tabcolsep6pt
\begin{tabular}{@{}lccc@{}}
\toprule
membership & $n$ & age (Gyr) & residual \\
\midrule
stage 1 (kinematic candidates, contaminated) & 259 & 1.89 & 0.93 \\
kinematic truth                              & 219 & 2.45 & 0.37 \\
two-stage (spectral rejection)               & 237 & \textbf{2.65} & \textbf{0.17} \\
\bottomrule
\end{tabular}
\caption{Isochrone ages for M~67 from three membership lists, against the
literature age of 2.82 Gyr. The two-stage pipeline --- kinematic candidates,
then rejection of spectral outliers --- recovers the age better than the raw
kinematic membership does.}
\label{tab:m67}
\end{table}

Read the middle row carefully, because it is the surprising one. The kinematic
membership is the \emph{truth} for every benchmark in this workbook, and it
still gets the age wrong by 0.37 Gyr. The two-stage list, which throws away 22
stars that kinematics accepted, gets it wrong by 0.17. The contaminants that
the spectral stage removed were not random: they were young field stars with
the right position and velocity, dragging the fitted age down.

Two honest qualifications. This is one cluster, so it is an existence proof
that membership quality propagates into parameters, not a calibration of the
size of the effect. And the effect is measured in the same direction as the
known bias of contamination on a young cluster's age --- which is reassuring,
and also means the mechanism is understood rather than merely observed.

\subsection{What this does and does not establish}

The Galactic-archaeology promise of chemical tagging is large: if the chemical
fingerprint of a birth cluster survives dispersal into the field, then the
field's stars can be assigned to birth sites, and the disc's star-formation
history can be reconstructed from the field alone, without the kinematic
luck that today's analyses depend on. This workbook measures where that
programme stands at our precision, and the honest summary is:

\begin{itemize}
\item \textbf{Established here.} Spectra contain cluster information that
  element-ratio summaries do not fully capture (\S\ref{sec:spectral}); that
  information is enough to improve membership lists well beyond kinematic
  selection in the specific case we tested (\textbf{Table~\ref{tab:m67}}), and
  it improves a physical parameter as a result; and distances from member
  giants agree with the literature for the clusters where we can measure them.
\item \textbf{Not established.} That the improvement generalises. One cluster,
  one pipeline, and a sample with duplicate rows and a dominant globular
  cluster. A referee will ask for the same experiment on the 25-cluster set
  with five independent membership definitions, and they will be right to.
\item \textbf{Not attempted here.} Recovering a dispersed cluster from the
  field with no kinematic prior. \textbf{Table~\ref{tab:field}} is our best
  attempt and it reaches precision $\approx 0.62$ at recall $0.42$ --- good
  enough to sharpen a membership list, not good enough to find a cluster from
  scratch. Strong chemical tagging remains an open problem in the sense that
  \citet{Casamiquela:21} define it: their best-case sample, with differential
  abundances and no field stars, recovers 9 of 31 clusters at the 40\,\%
  threshold and one at 70\,\%, and their clustering step recovers less than that
  on our stars (\S\ref{sec:literature}).
\end{itemize}

\begin{issues}[EXERCISES]
\begin{enumerate}
\item Reproduce the two-survey fit for one of the other clusters in the
  combined sample. Report age, distance and their posterior widths, with and
  without the red-clump prior. How much of the age improvement comes from the
  main sequence and how much from the prior?
\item The age--metallicity degeneracy means the posterior stays broad even
  with perfect membership. Choose a cluster in the sample and quantify the
  degeneracy directly: plot the posterior in the age--metallicity plane and
  measure its correlation coefficient. What extra observable would break it?
\item Contamination bias has a sign. Construct a synthetic cluster with 10%
  young field stars added, fit it with the same machinery, and compare the
  recovered age with the truth. Repeat with old contaminants. Which is the
  larger effect, and why?
\item \textbf{Table~\ref{tab:m67}} removes 22 stars from the kinematic list
  and improves the age. Inspect those 22 stars: what are their abundances,
  their kinematics and their position relative to M~67? Write the three-sentence
  case for and against calling them contaminants.
\end{enumerate}
\end{issues}

% 16 · The assignment
\section{THE ASSIGNMENT}
\label{sec:assignment}

\deckslides{slide 50, the hands-on brief.}

This chapter is the brief. Everything before it is the material you need to
execute it; everything you need to execute it is in this repository.

\subsection{What you are producing}

One object each (your assigned cluster) and one short report. The report makes
exactly three claims, each with evidence:

\begin{enumerate}
\item \textbf{A membership list} for your cluster, with the criterion stated
  precisely enough that another student can reproduce it: the cuts, the
  parameters, the random seed, the commit hash.
\item \textbf{An honest score} for that list, against both available referees
  (kinematic selection and Simbad memberships), reported as recall
  \emph{and} precision, and with the degenerate cases printed
  (\S\ref{sec:validation}).
\item \textbf{One physical parameter} derived from the list (an isochrone age
and distance, or a red-clump distance) with its uncertainty, and a sentence on
how the answer would change if your membership were 20\% larger or smaller.
\end{enumerate}

A report that claims a method is ``better'' without stating the population,
the task, the seed count and both halves of the metric pair has not met the
brief. This is the whole content of \S\ref{sec:validation}, applied to you.

\subsection{Setup}

The environment is reproducible from the repository root:

\begin{verbatim}
uv sync                     # python + dependencies
uv run cluster --help       # the CLI
uv run cluster download     # APOGEE DR19 + GALAH DR4 + Gaia DR3 subsets
\end{verbatim}

Field runs are the slow ones. For interactive work, set \texttt{CLUSTER\_FAST=1}
to cap the field sample; every number quoted in this workbook was produced with
\texttt{CLUSTER\_FAST=0}. A container with the same environment is provided
(\texttt{Dockerfile}), and the full walkthrough with worked examples is
\texttt{docs/student\_activities.md}.

\subsection{Your cluster}

\begin{table}[htbp]
\centering
\small
\tabcolsep4pt
\begin{tabular}{@{}llll@{}} \toprule cluster & kind & note & suggested track \\
\midrule M~3 & globular & metal-poor ($-1.6$) & A (spectral wins \\ M~5 &
globular & & A \\ M~13 & globular & & A \\ M~15 & globular & metal-poor
($-2.2$), southern & A or C \\ M~71 & globular & & A \\ M~107 & globular &
southern (DEC $-13^\circ$) & C \\ M~92 & globular & metal-poor ($-2.3$) & A \\
Pleiades & open & young ($0.1$ Gyr); chemistry fails here & A (expect a null)
\\ King~7 & open & sparse, faint & A \\ Berkeley~71 & open & sparse & A \\
IC~166 & open & sparse, contaminated & A \\ NGC~2158 & open & & A \\ NGC~1245 &
open & & A \\ King~5 & open & & A \\ NGC~7789 & open & giant-rich ($1.4$ Gyr) &
B) clump \\ NGC~1798 & open & & A \\ NGC~2420 & open & & A \\ NGC~6819 & open &
giant-rich; clump residual $+0.01$ & B \\ M~67 & open & the workhorse; most
data to beat & C (or A) \\ Berkeley~66 & open & sparse & A \\ NGC~188 & open &
old ($7.6$ Gyr) & B \\ NGC~6791 & open & old, metal-rich ($+0.4$) & B (or A) \\
Berkeley~17 & open & & A \\ NGC~2243 & open & the sweet spot: main sequence
\emph{and} clump & B or C \\ Collinder~261 & open & the sweet spot: 214 GALAH
members & C \\ \bottomrule
\end{tabular}
\caption{The 25 clusters and the track each one suggests. Track assignments
are advisory; the reasons are in the notes and in
\S\ref{sec:benchmark}--\S\ref{sec:physics}. Metadata for every cluster,
including the literature ages and distances you will be scored against, is in
\texttt{src/cluster/literature.py}.}
\label{tab:assignment}
\end{table}

Pair up on the open clusters; they are the hard part. Globulars tag cleanly
already, which makes them a faster and solo-friendly track, and it is worth
asking, while you work on one, why that is (\S\ref{sec:benchmark}: metallicity,
age and the absence of recent star formation all push in the same direction).

\subsection{The three tracks}

\textbf{Track A: spectral embeddings.} Compare the 16 abundances with the
256-dimensional spectral latent on your cluster:

\begin{verbatim}
uv run cluster run --spectral data/embeddings/attention_broad_merged.parquet
\end{verbatim}

Best on the metal-poor globulars, where element ratios lose discriminating
power and the spectrum retains it. Expect a null result on the Pleiades, and
report it as a result.

\textbf{Track B: isochrones and the red clump.} Recover age and distance:

\begin{verbatim}
uv run python scripts/red_clump.py
uv run python scripts/sweet_spot.py
\end{verbatim}

Best on the giant-rich open clusters (NGC~6819, NGC~7789). Compare against
\texttt{src/cluster/literature.py} and be explicit about the age--metallicity
degeneracy of \S 15.1.

\textbf{Track C: two surveys.} GALAH main sequence plus APOGEE giants:

\begin{verbatim}
uv run python scripts/build_galah_apogee.py
uv run python scripts/rerun_combined.py
\end{verbatim}

Southern clusters only (DEC $\lesssim +25^\circ$). This is the track with the
largest payoff and the largest number of ways to go quietly wrong: the parallax
relaxation is required, not optional, and the two surveys' abundance scales must
be reconciled before the combined fit means anything.

\subsection{Levers you may turn}

One axis at a time, and log the configuration each time: the sweep in
\textbf{Table~\ref{tab:levers}} was only interpretable because every run
recorded its settings.

\begin{table}[htbp]
\centering
\small
\tabcolsep5pt
\begin{tabular}{@{}lll@{}}
\toprule
lever & where & what it changes \\
\midrule
\texttt{TSNE.perplexity} & \texttt{src/cluster/config.py} & neighbourhood size of the embedding \\
\texttt{UMAP.n\_neighbors} & \texttt{src/cluster/config.py} & local graph scale \\
\texttt{HDBSCAN.min\_cluster\_size} & \texttt{src/cluster/config.py} & the size floor (removed by EVoC) \\
\texttt{SNR\_MIN} & environment & sample quality \\
\texttt{NORMALIZE\_ROWS} & environment & Euclidean $\leftrightarrow$ cosine equivalence \\
\texttt{USE\_ELEMENT\_WEIGHTS} & environment & element weighting, post-standardisation \\
\texttt{DWARF\_ONLY} & environment & giant/dwarf selection \\
\bottomrule
\end{tabular}
\caption{The levers, and the one rule that applies to all of them: change one
at a time, and report the effect with both halves of the metric pair. The
test-suite defaults are in \texttt{src/cluster/config.py}; the measured effects
on the DR17 sample are in \texttt{docs/experiment\_results.md}, and they do not
all survive on the current data.}
\label{tab:levers2}
\end{table}

\subsection{How you are assessed}

\begin{issues}[THE RUBRIC]
\textbf{Reproducibility (40\%).} A second person, on a clean checkout, gets
your numbers. Configuration, seed, commit hash and the exact command are in
your report. If you used a subsample, say how it was drawn.

\textbf{Empirical care (40\%).} Both referees. Both halves of recall and
precision. Mean $\pm$ s.d. over at least three seeds for any stochastic step.
The largest-group fraction printed. A null result reported as a null result: a
cluster that chemistry cannot separate is a finding, not a failure.

\textbf{Physics (20\%).} One parameter with an uncertainty, a comparison with
the literature value, and an honest statement of the dominant systematic, which
in most cases will be your membership list, not your photometry.
\end{issues}

\begin{issues}[EXERCISES]
\begin{enumerate}
\item Write the one-paragraph methods section of your report before you run
  anything: the population, the cuts, the method, the seed, the metrics, the
  referees. If you cannot write it, you are not ready to run.
\item Take a published tagging result in the field (a table of members, or a
claimed purity) and try to score it with this workbook's protocol. What
information is missing from the paper, and what would you have asked the
authors for?
\item Your cluster is listed with a suggested track. Write two sentences
  arguing for the other track, using only the cluster's properties (age,
  metallicity, declination, richness).
\end{enumerate}
\end{issues}

% 17 · Conclusions: the eight lessons, and what comes next
\section{CONCLUSIONS}
\label{sec:conclusions}

\deckslides{slides 48--49: the eight lessons and the take-away.}

This chapter collects the eight lessons the benchmark produced, each with the
number that supports it and the section that demonstrates the mechanism. They
are the same eight the lecture closes on; here they are expanded to the length
the evidence allows.

\subsection{Eight lessons from the benchmark}

\textbf{1 · All-sky collapses; cut a region.} One embedding of every clean star
in the survey buries each cluster in the field: the field is 357\,056 stars of
chemically smooth background, and 25 clusters of a few dozen members each are
not density peaks in that space. The published answer, and the one that works,
is to restrict the sky. \citet{Kos:17} used $40^\circ$ around the Pleiades, and
our region runs use a cone scaled to the cluster, $\max(3^\circ, 10\times$
diameter$)$. The cost of ignoring this is measurable: in the uniform all-sky
run the abundance-only t-SNE arm reaches recall $0.201$ at precision $0.240$
(\textbf{Table~\ref{tab:field}}), while the same pipeline in region mode on the
DR17 data reached recall $0.73$ at precision $0.47$
(\textbf{Table~\ref{tab:curse}}). Same methods, same clusters; the difference
is the contrast between the cluster and its immediate neighbourhood.

\textbf{2 · Precision is the hard part.} Read across the region-by-region sweep
and the macro precision is $0.28$ for t-SNE, $0.08$ for UMAP and $0.01$ for
EVoC (25 clusters, DR19). The field is chemically similar to the clusters (that
is what it means for a cluster to be dissolved in it) so a method that recovers
members also recovers their look-alikes. We can now put a number on the
mechanism that \citet{Casamiquela:21} report: they find that more than 70\,\%
of the groups in their best-case solution are statistical mixtures, and our own
cluster-only runs give $64$--$72\,\%$ of groups that recover no cluster at all
(\texttt{stat40} in \S\ref{sec:literature}). Their negative result is about
cluster-to-cluster separation and ours about field contamination, but both say
the same thing: at survey precision, the chemical differences between clusters
of the same age are comparable to the differences between a cluster and the
field.

\textbf{3 · Recall $\approx 1$ is usually a blob.} HDBSCAN$^*$ on the abundance
space swallows the dense field into a single cluster, after which any
best-overlap scorer reports that the true members are inside it: recall $1.000$
at precision $0.0013$ (UMAP on abundances, \textbf{Table~\ref{tab:field}}, and
note that this run has row-normalisation \emph{on} (it is the configuration
default, and it does not prevent the collapse). Two defences are built into
this workbook: report kNN purity, which cannot be gamed by a group's size; and
never quote recall without its precision. The third defence we tried)
row-normalisation. Is the subject of the quiet lesson below.

\textbf{4 · Globulars tag cleanly; open clusters do not.} The kNN purity of the
metal-poor globulars reaches $0.4$--$0.5$ while solar-metallicity open clusters
stay below $0.2$. Three effects push the same way: globulars are older, more
metal-poor, and formed from gas that no longer resembles the present-day
interstellar medium, so their chemistry is distinct from the field they sit in.
The Pleiades are the limiting case ($100$ Myr old, formed from gas the field
still resembles) and the prediction that chemistry fails there is part of the
assignment (\S\ref{sec:assignment}), to be reported as a null result when it
does.

\textbf{5 · t-SNE isolates, UMAP over-merges, EVoC is fast but coarse.} Three
sets of assumptions, three answers, and no winner. t-SNE produces the most
separated picture and is the least stable under row order
($0.218$--$0.560$ on the abundance matrix, \S\ref{sec:validation}). UMAP is
the fastest and the only arm we run on the full field, and it is the arm that
produces blobs. EVoC needs no clustering choice from the user, works on the
cosine geometry the data actually have, and has a seed spread
($\pm0.043$ on abundances) that overlaps the gaps being claimed between
methods.

\textbf{6 · Levers are fragile, and read them in pairs.} On the DR17 data,
weighting each element by $1/\sigma$ \emph{after} standardising lifted M~67's
t-SNE recall from $0.09$ to $0.44$: the largest single effect in the sweep, and
the fix for a bug that had made the same switch a no-op. On the DR19 data the
same switch moves recall from $0.10$ to $0.96$ while precision falls from
$0.14$ to $0.05$: t-SNE's cluster grows until it swallows the field. The region
sweep agrees, with precision falling $0.83 \rightarrow 0.24$. A lever selected
on one metric, one cluster and one seed is a hypothesis; it becomes a result
only when it is re-run on the current data with both halves of the metric pair
reported.

\textbf{7 · Spectra beat abundances, and not PCA.} A self-supervised masked
autoencoder on the raw spectrum separates the 25 clusters better than the 16
ASPCAP abundances ($0.760$ vs $0.578$ homogeneity, UMAP, the same 982 stars),
and a plain PCA of the same spectra does about as well ($0.754$). Two very
different representations of the same pixels land on the same score, which
tells you where the limit is: not in the model, but in what the spectra and the
metric jointly permit. On field retrieval the two tie to the third decimal
($0.418$/$0.619$ against $0.417$/$0.631$), and the only field-retrieval claim
we made for the spectral latent has been withdrawn for the reason given in
\S\ref{sec:provenance}.

\textbf{8 · Two surveys, two parameters.} APOGEE sees giants: the red clump
gives distances, and NGC~6819 comes out within $0.01$ mag of the literature
value (\textbf{Table~\ref{tab:dm}}). GALAH sees the main sequence and turnoff:
that is where age sensitivity lives. The joint fit is not there yet (the age
posterior ($\sigma_{\log a} \approx 0.97$) is close to the width of the flat
prior, so the age half of the promise needs a better likelihood before it shows
anything) but the distance half works, and the NGC~2243 experiment
(\S\ref{sec:physics}) shows what the combination can do when the distance is
pinned and the main sequence is added: an age residual improving from $0.52$ to
$0.09$.

\subsection{The quiet lesson}

We used to say that row-normalisation breaks the blob, with a precision gain
of three- to fivefold. Measured across the 25 clusters of the DR19 region sweep,
it does not: turning normalisation \emph{off} raises UMAP's mean precision by
$1.7\times$ ($0.08 \rightarrow 0.14$) and cuts the number of blobs from twelve to
five, while t-SNE is unchanged ($0.28 \rightarrow 0.26$) and EVoC moves by a
few thousandths. On M~67 the effect runs the other way: with normalisation on,
t-SNE precision is $0.83$; off, it is $0.28$. The lever is real, but its
direction depends on the method and the cluster, which is exactly what a
lever selected on one cluster at one seed looks like when it is re-measured
properly (\S\ref{sec:validation}). What survives as a rule is smaller and
duller: the precision--recall trade is there in every setting, and it is a
choice about which error you can live with, made explicit.

\begin{summary}[SUMMARY POINTS]
\begin{enumerate}
\item Clustering is a decision problem: which space, which metric, which notion
of ``cluster'', which scale, and every choice belongs in the write-up.
\item Geometry beats algorithm. The largest improvements in this project came
  from changing the representation (spectra instead of abundances,
  row-normalisation, 64 components instead of 256), not from changing the
  clusterer.
\item Density-based clustering needs a hierarchy or a persistence criterion.
  One global threshold cannot serve a field of 357\,056 stars and a cluster of
  6 members at the same time.
\item Embeddings are for neighbourhoods, not distances. A t-SNE or UMAP map is
  evidence about who is next to whom, and about nothing else.
\item Validate with something the algorithm never saw, and print the
  degenerate cases. Largest-group fraction, number of groups, both halves of
  the metric pair.
\item Measure the ceiling and report the gap. Kinematics are not a competitor;
  they are the ceiling that chemistry is being compared with.
\item A lever is a hypothesis. Re-run it on the current data, with both
  halves of the metric pair, before quoting it.
\item Membership quality propagates into physics: M~67's isochrone age moves
  from $1.89$ to $2.65$ Gyr, against a literature value of $2.82$, as the
  membership list improves.
\end{enumerate}
\end{summary}

\begin{issues}[FUTURE ISSUES]
\begin{enumerate}
\item \textbf{Re-embed on a single data product.} The mixed-product latent of
  \S\ref{sec:provenance} invalidates every field-retrieval number computed from
  it. Uniform \texttt{mwmStar} spectra for the backfilled members are already
  downloaded ($736$ spectra), and the re-embedding run awaits the model
  checkpoint. Until it is done, the field-retrieval rows for the spectral
  latent should be read as provisional.
\item \textbf{Collapse the duplicate rows.} 298 member rows share an
  \texttt{APOGEE\_ID} with another row. Until the sample is deduplicated, every
  macro average over clusters is weighted by something that is not a star.
\item \textbf{Benchmark against \citet{Spina:25}.} \S\ref{sec:literature}
covers the published negative result (their clustering step on our stars, their
metric triple on our scores) but a common-population comparison with the
graph-attention autoencoder, on the same clusters and the same two tasks, is
still the experiment this work is missing.
\item \textbf{Test the strong-tagging claim.} Everything here is measured with
  a kinematic prior available. Recovering a dispersed cluster from the field
  alone remains the open problem, and the honest current status is
  precision $\approx 0.62$ at recall $0.42$.
\item \textbf{A metric for high-dimensional chemical space.} The dilution
  result of \S\ref{sec:benchmark} is a statement about Euclidean distance, not
  about chemistry. A metric that respects the compositional structure of
  abundance vectors might make concatenation behave, and nobody has
  demonstrated one.
\end{enumerate}
\end{issues}

\begin{issues}[DISCLOSURE AND PROVENANCE]
Every number in this workbook is produced by code in the accompanying
repository at a recorded commit, and the scripts that generate the tables are
named in the sections that use them. Three things are worth distinguishing.
(1) The duplicate-row and row-order audit of \S\ref{sec:validation}, the
product-mismatch diagnosis of \S\ref{sec:provenance}, and the withdrawal of
the field-retrieval claim for the mixed latent are the repository's own
findings, reported here as they stand. (2) The K-means row of
\textbf{Table~\ref{tab:honest}} was added for this workbook
(\texttt{article/scripts/kmeans\_baseline.py}) and had not been measured
before. (3) The figures are the lecture's, reused from the deck so that the
workbook and the slides describe the same drawings; the animations are
flattened to filmstrips, and the deck itself is the primary reference for
anything that moves.
\end{issues}

\subsection{Acknowledgments}

The workbook accompanies the 2026 IAA--SO school on advanced neural networks,
and is built on the accompanying repository's benchmark. The data are the
public work of the SDSS-V/APOGEE, GALAH and Gaia collaborations, and the
methods are the work of the authors cited throughout. The teaching figures and
the dataset design come from the lecture; the isochrone fits use the PARSEC
grids through AStECA and \texttt{emcee}; the pipeline is written on
\texttt{numpy}, \texttt{pandas}, \texttt{scikit-learn}, \texttt{astropy},
\texttt{matplotlib} and \texttt{umap-learn}, with EVoC for the production
runs. Errors of interpretation that remain are ours, and the exercises exist
so that the next pass is yours.

\begin{figure}[htbp]\centering
  \begin{tabular}{@{}c@{\hskip14pt}c@{}}
    \fig[0.22]{qr_school.png} &
    \fig[0.22]{qr_library.png}
  \end{tabular}
  \caption{The school programme and the reference library. The deck itself,
with its animations, is at the URL on the title page and in the QR code of
\textbf{Figure~\ref{fig:qrcodes}}.}
  \label{fig:qrclosing}
\end{figure}


\bibliographystyle{ar-style2}
\bibliography{references}

\end{document}
